%% ============================================================
%%  jmlr-paper-main.tex
%%  HypatiaX: Hybrid Symbolic-Neural Framework for
%%  Extrapolation-Reliable Analytical Discovery
%%  MAIN PAPER — JMLR submission (extended, v3.0 benchmark)
%%  Author: Ruperto Pedro Bonet Chaple
%%  Built: April 2026
%%  Companion files:
%%    supp_routing_improvements.tex  — Supplementary A
%%    supp_benchmark_report.tex      — Supplementary B
%% ============================================================

\documentclass[twoside,11pt]{article}

%% ── JMLR house style ────────────────────────────────────────
\usepackage{jmlr2e}

%% ── Core maths & fonts ──────────────────────────────────────
%% NOTE: amsthm is deliberately NOT loaded. jmlr2e.sty already
%% defines proof/theorem/lemma/proposition/remark/corollary/
%% definition (see jmlr2e.sty "PROOF, THEOREM, and FRIENDS" section)
%% using plain-LaTeX \newtheorem/\newenvironment. Loading amsthm on
%% top of that redefines the kernel theorem-optional-argument
%% machinery and breaks jmlr2e's already-declared environments
%% (symptom: "No counter 'Definition' defined" on
%% \begin{definition}[Title]). Neither this file nor the
%% supplementary files use any amsthm-only feature (\theoremstyle,
%% \qedhere, \newtheorem*), so it isn't needed.
\usepackage{amsmath,amssymb,amsfonts}

\usepackage{bm}
\usepackage{microtype}
\microtypesetup{expansion=false}

%% ── Figures, tables, algorithms ─────────────────────────────
\usepackage{graphicx}
\usepackage{float}
%%adjustbox-removed
\usepackage{booktabs}
\usepackage{caption}
\usepackage{subcaption}
\usepackage{silence}
\WarningFilter{caption}{Unknown}
\usepackage{algorithm}
\usepackage{algpseudocode}
\usepackage{multirow}
\usepackage{threeparttable}
\usepackage{longtable}
\usepackage{tabularx}
\usepackage{makecell}
\usepackage{colortbl}
\usepackage{array}

%% ── Typography & misc ───────────────────────────────────────
\usepackage{enumitem}
\usepackage{hyperref}
\usepackage{xcolor}

\graphicspath{{figures/}{../figures/}{hypatiax/data/results/figures/}}
% Accept pdf, png, jpg in that priority order — figures_deploy produces .png
% when generate_figures.py writes .png only; pdflatex accepts all three.
\DeclareGraphicsExtensions{.pdf,.PDF,.png,.PNG,.jpg,.JPG}

%% ── Colours ─────────────────────────────────────────────────
\definecolor{passgreenBg}{RGB}{204,255,204}
\definecolor{warnAmberBg}{RGB}{255,235,180}
\definecolor{failredBg}{RGB}{255,204,204}

%% ── Placeholder for missing data ────────────────────────────
\newcommand{\TODO}{\textcolor{red}{[TODO]}}

%% ── Math macros ─────────────────────────────────────────────
\newcommand{\Rsq}{R^{2}}
\newcommand{\Rtrain}{R^{2}_{\mathrm{train}}}
\newcommand{\Rtest}{R^{2}_{\mathrm{test}}}
\newcommand{\Redge}{R^{2}_{\mathrm{edge}}}
\newcommand{\II}{\mathrm{II}}
\newcommand{\abs}[1]{\left|#1\right|}
\newcommand{\wLLM}{w_{\mathrm{LLM}}}
\newcommand{\wNN}{w_{\mathrm{NN}}}
\newcommand{\R}{\mathbb{R}}
\newcommand{\verified}[1]{\textcolor{blue!60!black}{#1}}

%% ── Theorem environments ────────────────────────────────────
%% theorem/definition/proposition/corollary/remark/lemma are ALL
%% already defined unconditionally by jmlr2e.sty (identical names,
%% shared "theorem" counter) — see jmlr2e.sty lines 388-394. They
%% must NOT be redeclared here; only result/finding/constraint are
%% genuinely new and not touched by jmlr2e.
\newtheorem{result}{Empirical Result}
\newtheorem{finding}{Empirical Finding}
\newtheorem{constraint}{System Constraint}
\newenvironment{demonstration}{\begin{trivlist}\item[\hskip\labelsep\textbf{Empirical Demonstration.}]}{\end{trivlist}}
\newenvironment{argument}{\begin{trivlist}\item[\hskip\labelsep\textbf{Theoretical Argument.}]}{\end{trivlist}}

% =============================================================================
\title{HypatiaX: A Hybrid Symbolic-Neural Framework for
  Extrapolation-Reliable Analytical Discovery}

\author{R.~P.~Bonet Chaple}

\date{2026}

% =============================================================================
\begin{document}
\raggedbottom

\jmlrheading{1}{2026}{1-99}{1/26}{4/26}{Bonet Chaple}
\ShortHeadings{HypatiaX: A Hybrid Symbolic-Neural Framework for Extrapolation-Reliable Analytical Discovery}{Bonet Chaple}
\firstpageno{1}

\title{HypatiaX: A Hybrid Symbolic-Neural Framework for
  Extrapolation-Reliable Analytical Discovery}

\author{Ruperto Pedro Bonet Chaple\\
  \texttt{ruperto.bonet@modelphysmat.com}\\
  Independent Researcher, London, United Kingdom
}

\editor{}

\maketitle

% ── Author note on numerical accuracy ────────────────────────────────────────
\begin{center}
\fbox{\parbox{0.88\textwidth}{%
\small\textbf{Note on runtime claims.}
An earlier draft of this paper asserted a ``73\,\% runtime reduction''
(i.e.\ a $3.7\times$ speedup of HypatiaX over the neural-network baseline).
That figure is \emph{not supported} by the benchmark data (HypatiaX v3.0).
The verified runtime comparison is presented in full in Section~\ref{sec:timing};
the honest claim is a $1.73\times$ speedup on LLM-routed cases only
($n=68$ out of 74).  All other quantitative claims in this paper have been
independently verified against the v3.0 statistical report and are reported
exactly as measured.}}
\end{center}
\vspace{0.5em}

% =============================================================================
\begin{abstract}
The hardest test of a formula-recovery system is not interpolation accuracy but
out-of-distribution extrapolation.  Two cases from our benchmark make this stark.
On Arrhenius, PySR-only collapses catastrophically at $5\times$ the
training range (far-$\Rsq = -12.56$); HypatiaX anchors the search via an
LLM warm-start and recovers the correct functional form (far-$\Rsq = 0.90$),
a genuine rescue rather than a shared failure. On Portfolio
Variance, PySR-only fails on all five tested random seeds; HypatiaX recovers
the exact closed-form formula ($\Rsq=1.000$ at $5\times$ training range) on
seeds 777 and 2024, and is close to or somewhat better than PySR-only on
seeds 42 and 123. We also document an honest failure mode on the opposite
side: on Michaelis-Menten, HypatiaX's extrapolation evaluation never
completed (no value recorded near, medium, or far), while PySR-only
succeeded (far-$\Rsq = 0.65$); and on Portfolio Variance seed~99, HypatiaX
underperforms PySR-only at far extrapolation. These phenomena---selective
rescue, selective degradation, and outright missing evaluations---are the
empirical core of this paper.

We introduce \textbf{HypatiaX}, a hybrid symbolic--neural framework that
addresses this through a five-stage routing and ensembling mechanism---trust
gating, transcendental complexity detection, neural degradation probing,
out-of-distribution detection, and uncertainty-weighted ensembling---that
dynamically selects between LLM symbolic output and a residual neural network
using training-data diagnostics alone.  The key insight is that the two paradigm
failure modes are complementary: routing between them eliminates catastrophic
failures while preserving the accuracy ceiling of the stronger component.

On a benchmark of 74 DeFi mathematical tasks under an aggressive
40\,\%/60\,\% PCA-directed extrapolation split, a standalone neural MLP achieves
median test $\Rsq = -0.47$, and a pure LLM baseline suffers 6 catastrophic
failures ($\Rsq < -10$) while achieving only 62.2\,\% near-perfect recovery.
Neither paradigm alone is sufficient.

Across all 74 tasks, HypatiaX achieves a median test $\Rsq$ of $1.00$ and a
near-perfect success rate ($\Rsq>0.99$) of \textbf{89.2\,\%}\footnote{This
figure is computed from the hybrid system's own reported \texttt{success}/$\Rsq$
fields without cross-checking the routing decision against the actual success
of the sub-method it names. A code-level attribution bug inflates it; the
corrected figure is 60.8\,\% (45/74), verified directly against
\texttt{hypatiax\_defi\_benchmark\_v3\_results\_seed42.json}. See
\S\ref{sec:hybrid-attribution-bug}.}---a $+27$
percentage point gain over the LLM baseline and $+83.8$\,pp over the neural
network---while \textbf{eliminating all catastrophic failures in HypatiaX's
own reported output}\footnote{The underlying computation is not free of
catastrophic failures: 22 of the 74 tasks route to a sub-method (typically
the LLM proposal) that itself failed catastrophically (e.g.\
\texttt{pure\_llm.test\_r2} $\approx -141{,}000$ on Loan-to-Value), but the
hybrid system's own accounting reports \texttt{success=True},
$\Rsq \approx 1.0$ regardless. The catastrophic failure is masked, not
eliminated. See \S\ref{sec:hybrid-attribution-bug}.}.
Gains scale with difficulty in the uncorrected figures: $+38.1$\,pp on hard
transcendental tasks (76.2\,\% vs.\ 38.1\,\% for the LLM alone)
\textbf{[flagged: corrected gain on Hard is $-4.8$\,pp, i.e.\ a loss, not a
gain; see \S\ref{sec:hybrid-attribution-bug}]}.
On the standard Nguyen-12 suite, HypatiaX achieves 11/12 recovery (91.7\,\%)
versus 10/12 for PySR-only and 0/12 for the neural baseline.
When the LLM formula is trusted (68 of 74 cases, 92\,\%), HypatiaX delivers a
$1.73\times$ median speedup over neural-network inference.
On 30 Feynman equations spanning mechanics to quantum physics, HypatiaX
achieves 12/30 near-perfect recovery under a random 80/20 split and 13/30
under a PCA-directed 40\,\%/60\,\% split\footnote{Both figures use a single,
strict $\Rsq \geq 0.999999$ threshold, applied identically to both splits;
near-misses ($\Rsq \in [0.999, 0.999999)$) are disclosed separately in
\S\ref{sec:feynman30} and never counted toward these figures. See the
threshold policy statement in \S\ref{sec:feynman30} for the single
convention used throughout this paper.}, versus 0/30 for the neural baseline
under both protocols (see Section~\ref{sec:feynman30}), establishing that reliability gains extend beyond the DeFi domain.
These results establish \emph{reliability under distribution shift}, not raw
accuracy, as the defining advantage of hybrid symbolic--neural systems.
\end{abstract}

\begin{quote}\small
\textbf{Note.} Several claims and one full table in this manuscript have
been withdrawn or flagged pending re-derivation during an internal review;
see \texttt{unconfirmed\_paper\_claims.tex} for the complete, consolidated
list with sourcing. Affected passages are marked in place.

\medskip
\noindent\textbf{Note (added, hybrid decision-attribution bug).} Every
near-perfect success-rate figure for the hybrid system quoted in this
manuscript (89.2\,\% overall, and the 100.0\,\%/89.7\,\%/76.2\,\% difficulty
breakdown) is computed directly from the pipeline's \texttt{hybrid.success}
/ \texttt{hybrid.test\_r2} fields \emph{without} cross-checking the
\texttt{hybrid.decision} field against the actual \texttt{success} flag of
the sub-method that decision names. Re-analysis of the underlying result
JSON (\texttt{hypatiax\_defi\_benchmark\_v3\_results\_seed42.json}) confirms
this is a real code-level bug, not evaluation noise: in 24 of 74 tasks,
\texttt{hybrid.decision} names a sub-method (always \texttt{"llm"} in this
run) whose own \texttt{success} field is \texttt{False} for that same task
(the method the hybrid record claims to have used had itself failed), yet
\texttt{hybrid.test\_r2} still reports a near-perfect ($>0.99$) value in 22
of those 24 cases. Cross-checked against seeds
42/99/123/777/2024 for the correlated-portfolio subset, the pattern is
deterministic (5/5 seeds), confirming a code defect rather than a flaky
result. See \S\ref{sec:hybrid-attribution-bug} for the full quantification
and corrected figures; all affected passages are marked in place with a
forward reference to that section.
\end{quote}


\tableofcontents
\newpage

% =============================================================================
\section{Introduction}
\label{sec:intro}

Extrapolation beyond the observed training distribution is a fundamental challenge
in machine learning.  In financial and decentralised-finance (DeFi) domains, the
problem is particularly acute: pricing models, risk metrics, and market-maker
invariants must generalise correctly to input regimes that may differ dramatically
from observed historical data.  A system that interpolates well but extrapolates
poorly is of limited practical value in this setting.

Two broad paradigms have been applied to mathematical function recovery in such
domains:

\paragraph{Neural network regression.}
Multilayer perceptrons and related architectures are powerful interpolators but
are known to fail systematically under distribution shift.  Even with augmentation
techniques, neural networks trained on $\mathcal{D}_{\mathrm{train}}$ do not, in
general, recover the functional form of $f$ beyond the training convex hull.  Our
benchmark confirms this: the neural MLP baseline achieves a median test $\Rsq$ of
$-0.47$ under extrapolation---worse than predicting the mean.

\paragraph{LLM-based symbolic regression.}
Large language models can, in principle, recover exact closed-form expressions from
task descriptions, achieving perfect extrapolation when the correct formula is
produced.  However, this approach exhibits a severe bimodal failure pattern: on
a benchmark of 74 DeFi tasks with $K=30$ independent runs per task, LLMs achieve
median test $\Rsq=1.00$ but mean $\Rsq=-0.76$, with 6 catastrophic failures
($\Rsq < -10$).  More importantly, a single-run evaluation---the standard in prior
work---conceals run-to-run instability: on complex transcendental tasks, the same
prompt yields radically different expressions across runs (Instability Index
$\II > 0.10$).

\medskip\noindent
\textbf{Central claim.}
Neither paradigm alone is sufficient.  The key insight of HypatiaX is that
these two failure modes are \emph{complementary}: neural networks fail on
the tasks where LLMs succeed (transcendental extrapolation), and LLMs fail
on the tasks where a neural network with LLM-augmented features can stabilise
performance.  A routing mechanism that dynamically selects between the two
modes---based on observable diagnostics computed from training data alone---can
eliminate catastrophic failures while preserving near-perfect accuracy on
the cases where symbolic recovery is reliable.

\subsection{Contributions}\label{subsec:contributions}
\label{sec:contributions}

This paper makes the following contributions:

\begin{enumerate}[leftmargin=2em,label=(\roman*)]
  \item \textbf{HypatiaX benchmark}: a corpus of 74 closed-form DeFi tasks
    with an aggressive PCA-directed extrapolation split (40\,\%/60\,\%), together
    with a unified evaluation pipeline (HypatiaX v3.0) that corrects four
    systematic biases in prior versions (data leakage, split misconfiguration,
    formula evaluation inconsistency, and weak trust gating).

  \item \textbf{HypatiaX hybrid system}: a five-stage routing and ensembling
    mechanism---trust gating, transcendental token detection, neural degradation
    probing, OOD detection, and uncertainty-weighted ensembling---that achieves
    89.2\,\% near-perfect success rate ($\Rsq>0.99$) across all 74 tasks
    \textbf{[flagged: hybrid decision-attribution bug inflates this figure;
    corrected value 60.8\,\%, see \S\ref{sec:hybrid-attribution-bug}]}.

  \item \textbf{Verified runtime analysis}: a rigorous decomposition of wall-clock
    time showing that the $1.73\times$ speedup of HypatiaX over the neural baseline
    applies specifically to the 68 LLM-routed cases, and that aggregate mean time
    (6.8\,s) is higher than the neural baseline (3.0\,s) due to fallback costs.
    This corrects a previously circulated 73\,\% reduction claim that was not
    supported by the v3.0 benchmark data.

  \item \textbf{Difficulty-stratified analysis}: a breakdown of performance across
    easy (24 tasks), medium (29 tasks), and hard (21 tasks) cases, demonstrating
    that HypatiaX's gains are largest and most practically important on the hardest
    transcendental and multi-asset formulas.
\end{enumerate}

% =============================================================================
\section{Related Work}
\label{sec:related}

\subsection{Symbolic Regression}\label{subsec:symbolic-regression}
\label{sec:related_sr}

Symbolic regression aims to discover explicit functional forms from data without
assuming a fixed parametric model.  Classical methods rely on genetic programming
\citep{koza1992genetic} and evolutionary search; the SRBench benchmark \citep{lacava2021contemporary}
provides a standardised multi-dataset evaluation suite.  The AI Feynman
system \citep{udrescu2020ai} demonstrated that dimensional analysis and neural
pre-screening can guide symbolic search toward physically meaningful equations.

A critical limitation of the symbolic regression literature is the evaluation
paradigm: systems are typically assessed on a single run or on the best expression
found across multiple runs.  This conceals run-to-run instability, which---as we
show---is the dominant failure mode of LLM-based symbolic regression on complex
formulas.  Our multi-run protocol (30 independent runs per task) and Instability
Index directly address this gap.

\subsection{LLMs for Mathematical and Scientific Reasoning}\label{subsec:llms-for-mathematical-and-scientific-rea}
\label{sec:related_llm}

Chain-of-thought prompting \citep{wei2022chain} and Minerva \citep{lewkowycz2022minerva}
demonstrate that LLMs can solve complex mathematical problems.  Codex \citep{chen2021evaluating}
showed strong code and formula generation capability.  More recent work has explored
LLMs as symbolic regression engines, generating formulas from natural-language
descriptions of datasets.

All of these works evaluate correctness in a single-shot or best-of-$K$ setting and
do not report variance across runs as a first-class metric.  Our benchmark is the
first, to our knowledge, to quantify LLM instability systematically across a large
corpus of financial formulas under repeated stochastic inference.

\subsection{Equation Learners, Differentiable SR, and Uncertainty-Aware Methods}\label{subsec:equation-learners-differentiable-sr-and-}
\label{sec:related_eq}

Several strands of recent work are directly relevant to the hybrid design in
this paper.  The Equation Learner (EQL; \citealt{sahoo2018eql}) embeds
differentiable operator networks with $L_0$ regularisation to produce sparse,
interpretable expressions, demonstrating that structural bias can be enforced
without explicit symbolic search.  DISCOVER \citep{liu2022discover} reformulates
symbolic regression as differentiable optimisation over a symbolic search space,
enabling gradient-based recovery of functional form.  Bayesian and
Gaussian-process approaches to symbolic regression \citep{jin2020bayesian}
provide principled uncertainty estimates alongside recovered expressions ---
a goal shared by the confidence-weighted ensemble in \textsc{HypatiaX}'s
routing cascade.  \textsc{HypatiaX} differs from these works by treating
\emph{extrapolation reliability} as a first-class evaluation criterion
alongside recovery rate, and by combining evolutionary symbolic search with an
explicit multi-layer validation cascade that filters expressions before
acceptance.

\subsection{Hybrid Symbolic--Neural Systems}\label{subsec:hybrid-symbolicneural-systems}
\label{sec:related_hybrid}

The idea of combining symbolic structure with neural learning has a rich history.
\citet{cranmer2020lagrangian} use neural networks to discover Lagrangians, then
fit symbolic expressions to the learned dynamics.  \citet{udrescu2020ai} alternate
between neural pre-screening and symbolic search.  PySR \citep{cranmer2023pysr}
provides an efficient evolutionary-symbolic regression framework that can be
combined with neural features.

Our approach differs in motivation: HypatiaX is designed not to improve accuracy
over either pure method in isolation, but to \emph{eliminate catastrophic failures}
and achieve reliability across the full difficulty spectrum of DeFi tasks.

\subsection{Extrapolation in Neural Networks}\label{subsec:extrapolation-in-neural-networks}
\label{sec:related_extrap}

The failure of neural networks to extrapolate is well-documented \citep{xu2021neural}.
Data augmentation with scaled inputs \citep{lim2021temporal} can improve extrapolation
marginally but does not address the fundamental limitation.  Our results confirm
that even a residual MLP with LayerNorm, SiLU activations, and scaled augmentation
achieves median $\Rsq=-0.47$ under our 40\,\%/60\,\% extrapolation split.

% =============================================================================
% === SECTION: Empirical Evidence (from jmlr_paper_final.tex §3) ===
\section{Empirical Evidence of LLM Performance Across Domains}
\label{sec:llm_limitations}

To evaluate whether large language models (LLMs) can perform analytical formula discovery without symbolic constraints or external solvers, we conducted a systematic study of pure LLM-based formula generation \citep{koza1992genetic,kambhampati2024llms}. No symbolic regression, constraint enforcement, or post-generation correction was applied.

\subsection{Experimental Design}\label{subsec:experimental-design}
\label{sec:llm_exp_design}

We evaluated GPT-4 \citep{openai2023gpt4} on 40 interpolation-based formula discovery tasks spanning five domains: classical physics (Kinetic Energy, Gravitational Force, Ideal Gas Law), chemistry (Arrhenius, Henderson-Hasselbalch, Rate Law), biology (Allometric Scaling, Michaelis-Menten, Logistic Growth), and decentralized finance (Impermanent Loss, Price Impact, Portfolio Variance, Value-at-Risk, Liquidation Price). Tasks were presented as structured prompts containing variable names, units, and sample data. Outputs were evaluated by executing the generated formula symbolically and computing $\Rsq$ on a held-out test set within the training range.

\subsection{Baseline Results: Bimodal Performance}\label{subsec:baseline-results-bimodal-performance}
\label{sec:llm_baseline}

Pure LLM performance is strongly bimodal. On well-known physical laws (Kinetic Energy, Ideal Gas Law, Gravitational Force), the LLM achieves $\Rsq = 1.000$ with exact symbolic recovery. On domain-specific and specialized equations---particularly in chemistry (Henderson-Hasselbalch: $\Rsq = -3.2$), biology (Michaelis-Menten: $\Rsq = -68.6$), and DeFi (Portfolio Variance: $\Rsq = -21.0$)---performance collapses catastrophically. LLM success correlates with estimated equation prevalence in web corpora (Spearman $\rho = 0.89$, $p < 0.001$), consistent with the hypothesis that pretraining data distribution influences recovery performance \citep{kambhampati2024llms}.

\subsection{Failure Mode Taxonomy}\label{subsec:failure-mode-taxonomy}
\label{sec:llm_failures}

We identify five failure modes, ordered by severity:
\begin{enumerate}
  \item \textbf{Silent semantic errors} (High): Structurally plausible formula with incorrect constants or signs; high in-range $\Rsq$ but catastrophic extrapolation.
  \item \textbf{Distributional reasoning failures} (High): Misuse of statistical functions (e.g.\ wrong sign for VaR z-score).
  \item \textbf{Unit/scale inconsistency} (Medium): Missing unit conversions or scale factors.
  \item \textbf{Non-executable output} (Medium): Pseudocode or invalid Python; cannot be evaluated.
  \item \textbf{Incomplete construction} (Medium): Partial derivations missing key terms.
\end{enumerate}

\subsection{Case Studies}\label{subsec:case-studies}
\label{sec:llm_cases}

\paragraph{Arrhenius equation.}
The Arrhenius rate constant is $k = A \exp\!(-E_a/RT)$. The LLM correctly identifies the exponential structure but systematically misestimates the activation energy $E_a$ by a constant factor, producing train $\Rsq = 0.997$ but extrapolation $\Rsq = -12.6$ at $5\times$ the training range. This is the archetypal \emph{scale-incompatibility failure}: the formula is structurally correct but parametrically wrong in a way that only manifests under distribution shift.

\paragraph{Michaelis-Menten.}
$v = V_{\max}[S]/(K_m + [S])$. The LLM recovers the functional form on some runs ($\Rsq = 1.000$) but collapses on others ($\Rsq < -68$), producing high run-to-run instability ($\II = 0.31$). This motivates the Instability Index as a diagnostic.

\subsection{Success Pattern Analysis}\label{subsec:success-pattern-analysis}
\label{sec:llm_success}

LLM guidance succeeds when: (1) the equation is well-represented in web text; (2) the functional form is algebraic rather than transcendental; (3) the domain is classical physics or standard chemistry. It fails systematically on post-2020 DeFi equations, multi-asset risk formulas, and transcendental biology equations---precisely the cases where symbolic regression provides the most value.

\subsection{Extrapolation Testing}\label{subsec:extrapolation-testing}
\label{sec:llm_extrap}

Even when the LLM achieves near-perfect interpolation ($\Rsq \approx 1$), extrapolation performance is unreliable. Across 40 tests, 6 equations suffered catastrophic extrapolation failure ($\Rsq < -10$). This motivates the hybrid architecture: the LLM provides structural guidance while symbolic regression corrects parametric errors and the neural network provides a stable fallback.

% =============================================================================
% === SECTION: Theoretical Framework (from jmlr_paper_final.tex §4) ===
\section{Theoretical Framework}
\label{sec:theory}

\subsection{Formal Definitions}\label{subsec:formal-definitions}
\label{sec:theory_defs}

\begin{definition}[Analytical Expression Discovery]
Let $\mathcal{D} = \{(x_i, y_i)\}_{i=1}^n$ where $x_i \in \mathbb{R}^d$, $y_i \in \mathbb{R}$~\citep{udrescu2020ai}. An analytical expression discovery system produces $E: \mathbb{R}^d \to \mathbb{R}$ satisfying:
\begin{enumerate}
  \item \textit{Symbolic form}: $E$ is expressible using operators $\mathcal{O}$ in closed form
  \item \textit{Physical consistency}: $E$ satisfies domain constraints $\mathcal{C}$
  \item \textit{Empirical accuracy}: Loss $\mathcal{L}(E, \mathcal{D}) < \epsilon$
  \item \textit{Extrapolation}: $E$ generalizes beyond training regime
  \item \textit{Parsimony}: $E$ has minimal complexity given above constraints
\end{enumerate}
\end{definition}

\begin{definition}[Interpolation--Extrapolation Decoupling]
\label{def:decoupling}
A system exhibits interpolation--extrapolation decoupling if $\exists$ tasks where $\Rtrain \approx 1$ but $\Rtest \ll 1$ on an out-of-distribution test set.
\end{definition}

\begin{definition}[Reproductive Behaviour]
\label{def:reproductive}
A generative model $G$ exhibits reproductive behaviour if its output distribution concentrates near its training manifold: $P(G(x) \notin \mathcal{M}_{\text{train}}) < \delta$ for small $\delta$, where $\mathcal{M}_{\text{train}}$ is the training data manifold.
\end{definition}

\subsection{Main Theoretical Result}\label{subsec:main-theoretical-result}
\label{sec:theory_main}

\begin{finding}[Interpolation--Extrapolation Decoupling]
\label{prop:interp_extrap_decoupling}
Neural networks trained on in-distribution data can achieve near-perfect interpolation ($\Rtrain \approx 1$) while exhibiting catastrophic extrapolation failure. Symbolic methods that recover the exact functional form are bounded by numerical precision alone.
\end{finding}

\begin{proof}[Proof sketch]
Neural networks achieve high in-distribution $\Rsq$ via local approximations (universal approximation~\citep{hornik1989multilayer}). These approximations are not constrained to match the true functional form; extrapolation requires the correct form. Symbolic methods that recover the exact $f^*$ achieve $\Rtest = 1 - O(\epsilon_{\text{num}})$ by construction.
\end{proof}

\begin{finding}[LLM Reproductive Behaviour]
\label{thm:llm_reproductive}
Autoregressive LLMs trained by maximum-likelihood exhibit reproductive behaviour: output probability concentrates near training-corpus examples. For analytical expression discovery, this implies performance degrades on expressions underrepresented in pretraining data.
\end{finding}

\begin{proof}[Proof sketch]
Maximum-likelihood training minimises $-\mathbb{E}[\log P_\theta(x)]$, concentrating probability mass near training examples. For rare equations (low corpus frequency), the model assigns low probability to the correct expression, increasing the risk of plausible-but-wrong outputs.
\end{proof}

\subsection{Implications for Discovery}\label{subsec:implications-for-discovery}
\label{sec:theory_implications}

Finding~\ref{thm:llm_reproductive} explains the domain-dependence pattern in Section~\ref{sec:llm_limitations}: LLM performance correlates with equation prevalence ($\rho = 0.89$) because reproductive behaviour ties output quality to training exposure. This motivates hybrid architectures that use LLM structural guidance while correcting via symbolic search.

\subsection{Functional Form Recovery and Inductive Bias Alignment}\label{subsec:functional-form-recovery-and-inductive-b}
\label{sec:theory_recovery}

Symbolic regression performs evolutionary exploration of expression space guided by fitness evaluation, without requiring proximity to textual training examples. For problems admitting closed-form representations, restricting search to symbolic structures improves extrapolation reliability by aligning the inductive bias of the model with the structure of the target function. By the No Free Lunch theorem \citep{wolpert1997no}, performance depends on alignment between hypothesis class and problem structure. When the true relationship admits a closed-form representation, this alignment holds for symbolic search and fails for purely neural approaches.

% =============================================================================
\section{Problem Formulation}
\label{sec:problem}

\begin{definition}[Symbolic Regression Task]
A task $i$ is defined by a dataset $\mathcal{D}_i = \{(x^{(j)}, y^{(j)})\}_{j=1}^{N}$
with ground-truth function $f_i: \mathbb{R}^{d_i} \to \mathbb{R}$ such that
$y^{(j)} = f_i(x^{(j)})$.  The objective is to recover a symbolic expression
$\hat{f}_i$ such that $\hat{f}_i \approx f_i$ on held-out test data.
\end{definition}

\begin{definition}[Extrapolation Setting]
Let $\mathcal{D}^{\mathrm{train}}_i$ and $\mathcal{D}^{\mathrm{test}}_i$ be a
partition of $\mathcal{D}_i$ such that the test inputs lie beyond the training range
along the principal direction of input variation:
\begin{equation}
  X^{\mathrm{test}} \notin [\min(X^{\mathrm{train}}),\, \max(X^{\mathrm{train}})]
  \quad \text{(approximately, along the first PC)}.
  \label{eq:extrap}
\end{equation}
Evaluation on $\mathcal{D}^{\mathrm{test}}_i$ under condition \eqref{eq:extrap}
is called \emph{extrapolation evaluation}.
\end{definition}

\begin{definition}[Coefficient of Determination]
For a predictor $\hat{f}$ evaluated on a set $\mathcal{S}$:
\begin{equation}
  \Rsq(\hat{f}, \mathcal{S}) = 1 - \frac{\sum_{(x,y)\in\mathcal{S}}(y - \hat{f}(x))^2}
                                         {\sum_{(x,y)\in\mathcal{S}}(y - \bar{y}_{\mathcal{S}})^2}.
  \label{eq:r2}
\end{equation}
For aggregate statistics, $\Rsq$ values are clipped to $[-10, 1]$ to prevent
extreme failures from dominating mean computations.  All reported aggregates use a
\emph{fixed denominator} of 74 (all standard cases), with NaN and failure cases
counted as 0.
\end{definition}

\begin{definition}[Extrapolation Gap and Stability Score]
\begin{align}
  \text{Extrapolation Gap}_i &= \Rtrain - \Rtest, \label{eq:gap}\\
  \text{Stability Score}_i   &= \Rtest / \Rtrain. \label{eq:stability}
\end{align}
A large gap indicates that the model fits the training distribution but fails
to generalise; a stability score near 1 indicates consistent performance across
the distribution shift.
\end{definition}

\begin{definition}[Catastrophic Failure]
A run is classified as a catastrophic failure if $\Rsq < -10$, indicating that
the predicted function is worse than predicting the mean by a factor of more
than 10.  Such failures typically arise from evaluation errors, division by
zero, or radically incorrect symbolic expressions.
\end{definition}

% =============================================================================
\section{Benchmark Design}
\label{sec:benchmark}

\subsection{Overview and Scope}\label{subsec:overview-and-scope}
\label{sec:bench_overview}

The HypatiaX DeFi Benchmark comprises \textbf{74 test cases} spanning the core
mathematical domains of modern decentralised finance.  All tasks are tractable,
admitting a closed-form or well-defined functional representation.  Ground-truth
outputs are computed analytically to machine precision, ensuring that any
deviation in $\Rsq$ reflects a model error rather than label noise.

Tasks are drawn from five domain categories:
\begin{itemize}
  \item \textbf{Automated Market Makers (AMMs)}: constant-product invariants,
    spot pricing, Uniswap V3 concentrated liquidity, impermanent loss calculations,
    AMM output amounts, price slippage, and liquidity-provider fee earnings.
  \item \textbf{Lending Protocols}: Loan-to-Value ratio, health factor, collateral ratio,
    liquidation prices for both long and short positions, leveraged-position notional,
    multi-collateral LTV, cross-margin available balance, reserve ratio, utilisation rate,
    borrow APY from utilisation, supply APY from borrow, and partial liquidation amount.
  \item \textbf{Derivatives Pricing}: Black--Scholes call and put prices,
    all four first-order Greeks (Delta, Gamma, Theta, Vega), put--call parity,
    forward price for derivative, call option intrinsic value, simple options
    moneyness, and convexity adjustment.
  \item \textbf{Risk and Portfolio Metrics}: Value at Risk at 95\,\% and 99\,\%,
    Expected Shortfall at 95\,\% and 99\,\%, multi-day VaR and ES scaling,
    incremental VaR, correlated portfolio VaR, component ES, portfolio VaR for
    two correlated assets, annualised portfolio tracking error, portfolio Sharpe
    ratio, and position margin ratio.
  \item \textbf{Staking and Protocol Mechanisms}: simple staking APY, compounding
    staking returns, staking reward for fixed lock-up, validator commission adjusted,
    slashing penalty, protocol reserve accumulation, funding rate cost, and simple
    staking APY.
\end{itemize}

\subsection{Difficulty Classification}\label{subsec:difficulty-classification}
\label{sec:bench_difficulty}

Tasks are categorised by difficulty into three tiers based on functional structure:

\begin{itemize}
  \item \textbf{Easy} ($n=24$): linear or simple rational relationships, such as
    Loan-to-Value, spot price from AMM reserve, collateral ratio, and simple staking APY.
    These tasks admit algebraic solutions with AST depth $\le 3$.

  \item \textbf{Medium} ($n=29$): nonlinear algebraic or exponential relationships,
    such as impermanent loss, Uniswap V3 liquidity range, constant-product price impact,
    and compounding staking returns.  These tasks require multi-step composition
    of algebraic operations.

  \item \textbf{Hard} ($n=21$): transcendental, probabilistic, or multivariate
    formulations.  This category includes Black--Scholes pricing (which requires
    the Gaussian CDF $\Phi$), all four Greeks, Portfolio Expected Shortfall for
    correlated assets, and multi-day VaR scaling.  These tasks are the primary
    source of LLM stochastic instability.
\end{itemize}

\begin{remark}
An earlier version of this benchmark comprised 71 tasks (Easy=24, Medium=27,
Hard=20).  The current benchmark (v3.0, 74 tasks) adds three additional hard
cases in multi-asset risk: correlated portfolio VaR, component ES, and
multi-collateral LTV.  All results in this paper refer to the 74-case v3.0 benchmark.
\end{remark}

\subsection{Data Generation}\label{subsec:data-generation}
\label{sec:bench_datagen}

For each task $i$, we generate $N=200$ samples using domain-specific input
simulators.  Input ranges are chosen to reflect realistic DeFi parameter regimes:
\begin{itemize}
  \item Strike prices $K \in [80, 120]$, spot prices $S \in [50, 150]$.
  \item Volatilities $\sigma \in [0.05, 0.80]$, risk-free rates $r \in [0.01, 0.10]$.
  \item Collateral ratios $C/D \in [1.1, 5.0]$, utilisation rates $u \in [0.01, 0.99]$.
  \item AMM reserves $(x, y)$ chosen such that $x \cdot y = k$ for constant $k$.
\end{itemize}
All inputs are sampled uniformly within their respective ranges to provide dense
coverage of the training domain.

\subsection{Extrapolation Protocol}\label{subsec:extrapolation-protocol}
\label{sec:split}

To rigorously assess out-of-distribution generalisation, we employ a
\textbf{PCA-directed aggressive extrapolation split}:

\begin{enumerate}[leftmargin=2em,label=\arabic*.]
  \item For each task $i$, compute the first principal component direction
    $v_1 \in \mathbb{R}^{d_i}$ of the input matrix $X \in \mathbb{R}^{N \times d_i}$.
  \item Project all $N=200$ samples onto $v_1$, obtaining scalar projections
    $z^{(j)} = \langle x^{(j)}, v_1 \rangle$.
  \item Sort samples by $z^{(j)}$ in ascending order.
  \item Assign the lowest 40\,\% of samples (by $z^{(j)}$) to the training set
    $\mathcal{D}^{\mathrm{train}}_i$ and the highest 60\,\% to the test set
    $\mathcal{D}^{\mathrm{test}}_i$, with a small overlapping band in the
    quantile transition region to prevent abrupt boundary effects.
\end{enumerate}

This protocol has two important properties.  First, unlike axis-aligned splits,
PCA-directed ordering ensures that extrapolation occurs along the dominant
direction of variation in multivariate settings, reflecting the most
challenging direction of distributional shift.  Second, the 40\,\%/60\,\%
partition is substantially more aggressive than the 80\,\%/20\,\% train/test
splits used in standard symbolic regression benchmarks, making this a stringent
test of extrapolation capability.

\begin{remark}[On the split design]
The partial overlap in the quantile transition region is a deliberate design choice
to avoid the ``cliff edge'' effect where a perfectly correct formula appears to
fail due to numerical boundary issues.  This overlap is small ($\sim$5\,\% of
samples) and does not substantially reduce the extrapolation challenge.
\end{remark}

\subsection{Pipeline Corrections in v3.0}\label{subsec:pipeline-corrections-in-v30}
\label{sec:r2_bugfix}
\label{sec:corrections}

HypatiaX v3.0 incorporates four corrections over prior versions that systematically
biased reported results:

\begin{enumerate}[leftmargin=2em,label=\arabic*.]
  \item \textbf{Data leakage in ensembling}: Earlier versions computed ensemble
    weights using test-set residuals.  v3.0 computes all ensemble weights from
    training residuals only, preventing any form of test information from
    influencing predictions.

  \item \textbf{Incorrect extrapolation splitting}: Prior versions used axis-aligned
    min/max splits, which can inadvertently include extrapolation samples in the
    training set for multivariate tasks.  v3.0 uses the PCA-directed ordering
    described above.

  \item \textbf{Inconsistent formula evaluation}: Earlier
    versions used multiple evaluation backends (SymPy, eval, NumExpr) without a
    unified type-coercion layer, causing silent mismatches.  A specific bug in
    \texttt{evaluate\_llm\_formula} used a hardcoded absolute sum-of-squares
    threshold (\texttt{ss\_tot > 1e-10}) that returned $\Rsq = 0$ for equations
    with outputs far below unity (e.g.\ gravitational forces $\sim 10^{-11}$\,N).
    v3.0 replaces this with a relative threshold and uses a single unified formula
    executor with consistent \texttt{float64} broadcasting semantics throughout.

  \item \textbf{Weak symbolic trust criteria}: Prior versions accepted LLM formulas
    with training $\Rsq > 0.1$.  v3.0 raises this threshold to $\Rsq > 0.5$ and
    adds pathological pattern detection (division by zero, NaN outputs, infinite
    predictions).
\end{enumerate}

These corrections are not cosmetic: they collectively account for a substantial
fraction of the performance gap between the ``true'' hybrid method and the inflated
numbers reported in earlier drafts.

% =============================================================================
\section{Methodology}
\label{sec:method}

\subsection{Architecture Overview}\label{subsec:architecture-overview}
\label{sec:method_overview}

HypatiaX combines three computational components---an LLM symbolic generator,
a residual MLP, and a routing/ensembling layer---as illustrated in
Figure~\ref{fig:architecture}.  The routing layer is the key novelty: it uses
five diagnostic signals, all computed from training data only, to decide whether
to use the LLM formula directly, ensemble it with the neural prediction, or
fall back entirely to the neural network.

\begin{figure}[h]
\centering
\includegraphics[width=\linewidth]{hypatiax_three_systems}
\caption{\textbf{HypatiaX architecture: three internal modules plus a
router.} The \emph{PySR module} (symbolic regression) and the \emph{LLM
module} (symbolic prior) each propose a candidate formula from the
training data $(X,y)$; the \emph{Hybrid Fusion module} combines them. The
Router/Stability Check is a separate, fourth component that selects one
of four outcomes based on training-set diagnostics only (no test
information is used at decision time): accept the PySR fit directly, fall
back on low PySR confidence, accept the LLM prior directly, or take the
combined hybrid path --- returning the best symbolic expression and its
$R^2$. \textbf{Note on terminology:} this figure's underlying graphic
labels these three modules ``System~1'', ``System~2'', and ``System~3''
respectively. That local numbering is specific to this figure only and
must not be confused with the paper's System~1/2/3 \emph{ablation
variants} (Section~\ref{sec:arch_variants}: System~1 = Hybrid~v50\_2, the
full pipeline; System~2 = symbolic-only with validation; System~3 = LLM
with symbolic fallback), which is the numbering used in
Figure~\ref{fig:routing_cascade}, Table~\ref{tab:five_systems_full}, and
throughout the rest of the paper. To avoid the collision, this caption
refers to the three modules above by name rather than by number.}
\label{fig:architecture}
\end{figure}

\subsection{Component 1: LLM-Based Symbolic Generation}\label{subsec:component-1-llmbased-symbolic-generation}
\label{sec:llm_gen}

Given a task description $\mathcal{T}_i$ (a natural-language specification of the
functional relationship, including variable names, units, and domain context),
the LLM produces a candidate symbolic expression:
\begin{equation}
  \hat{f}^{\mathrm{LLM}}_i \leftarrow \mathrm{LLM}(\mathcal{T}_i).
  \label{eq:llm_gen}
\end{equation}
The generated expression is parsed by the unified formula executor and evaluated
on both the training and edge subsets of the training data (the upper 20\,\%
of training samples by projection onto $v_1$, i.e.\ the samples closest to the
extrapolation boundary).

We define the training fit:
\begin{equation}
  {R^{2,\mathrm{LLM}}_{\mathrm{train}}} = \Rsq\!\left(\hat{f}^{\mathrm{LLM}}_i,\,
  \mathcal{D}^{\mathrm{train}}_i\right),
  \label{eq:r2_train_llm}
\end{equation}
and the edge fit:
\begin{equation}
  {R^{2,\mathrm{LLM}}_{\mathrm{edge}}} = \Rsq\!\left(\hat{f}^{\mathrm{LLM}}_i,\,
  \mathcal{D}^{\mathrm{edge}}_i\right),
  \label{eq:r2_edge_llm}
\end{equation}
where $\mathcal{D}^{\mathrm{edge}}_i$ is the upper-20\,\% subset of the training
data described above.

\subsection{Component 2: Neural Network Estimator}\label{subsec:component-2-neural-network-estimator}
\label{sec:nn}

The neural baseline is a \textbf{residual MLP} with the following architecture:
\begin{itemize}
  \item Input: $\mathbb{R}^{d_i}$ (task-specific input dimension).
  \item Hidden layers: 3 residual blocks, each consisting of Linear $\to$
    LayerNorm $\to$ SiLU $\to$ Linear, with a skip connection.
  \item Output: $\mathbb{R}$ (scalar prediction).
  \item Optimiser: Adam, learning rate $10^{-3}$, weight decay $10^{-4}$.
  \item Training data augmentation: $(X, y) \cup (1.2X, 1.2y)$, i.e.\ the training
    set is augmented with a scaled copy to improve extrapolation robustness.
  \item Training time limit: 120 seconds wall-clock per case, to ensure
    fair comparison with the LLM method (which does not require training time).
\end{itemize}

When HypatiaX routes a task to the neural fallback, the LLM prediction
$\hat{f}^{\mathrm{LLM}}_i(x)$ is appended to the input as an additional feature:
\begin{equation}
  x' = [x;\, \hat{f}^{\mathrm{LLM}}_i(x)] \in \mathbb{R}^{d_i + 1}.
  \label{eq:augmented_input}
\end{equation}
This \emph{residual learning augmentation} allows the neural network to learn
corrections to an imperfect symbolic formula rather than learning the function
from scratch, which substantially reduces the extrapolation burden when the LLM
formula captures the correct functional form but has parameter errors.

\subsection{Component 3: Five-Stage Routing and Ensembling}\label{subsec:component-3-fivestage-routing-and-ensemb}
\label{sec:routing}

\begin{figure}[H]
\centering
\includegraphics[width=0.82\linewidth]{hypatiax_algorithm1_routing_cascade_v2}
\caption{\textbf{Algorithm~1 --- HypatiaX routing cascade.}
Coloured regions indicate the originating system, using the System~1/2/3
ablation-variant numbering defined in Section~\ref{sec:arch_variants}
(\emph{not} the module names used in Figure~\ref{fig:architecture}, which
is a different, figure-local numbering --- see that figure's caption):
amber marks System~1 (Hybrid~v50\_2, the full pipeline --- here, its
extrapolation-detection and warm-restart phases); green marks System~2
(symbolic-only with validation, i.e.\ the cold-start PySR phase); purple
marks System~3 (LLM with symbolic fallback, i.e.\ the LLM query phase).
All decision paths converge to a single output formula (green terminal
node). The convergence lock badge confirms Proposition~\ref{prop:convergence}.}
\label{fig:routing_cascade}
\end{figure}

\begin{proposition}[Routing Cascade Convergence]
\label{prop:convergence}
Let $h^*_1 = \arg\max_{h \in \mathrm{HOF}_1} R^2(h, D)$ and
$h^*_2 = \arg\max_{h \in \mathrm{HOF}_2} R^2(h, D)$.
If the scale-compatibility guard (Stage~3) does not reject the LLM
prior $c^*$, or $S$ is non-empty after suppression, then
$$\mathbb{E}\bigl[R^2(h^*_2)\bigr] \geq \mathbb{E}\bigl[R^2(h^*_1)\bigr].$$
\end{proposition}

\begin{proof}[Proof sketch]
Phase~2 initialises PySR with $\mathrm{HOF}_1 \cup S \supseteq \mathrm{HOF}_1$.
Since evolutionary search is a monotone improvement operator over any
initial population dominating the current Pareto front, the warm-start
population is a strict superset of the cold-start population.
By the Pareto dominance argument, $R^2(h^*_2) \geq R^2(h^*_1)$
for a fixed random seed; in expectation over seeds, the inequality is
preserved when $S$ contains at least one expression structurally
distinct from $\mathrm{HOF}_1$.
\end{proof}

The routing logic implements five diagnostic stages in sequence.

\paragraph{Stage 1 --- Trust Gate.}
Accept the LLM formula only if it passes two criteria:
\begin{equation}
  \text{Trust}_i^{\mathrm{LLM}} = \mathbf{1}\!\left[{R^{2,\mathrm{LLM}}_{\mathrm{train}}} \ge 0.5\right]
  \;\land\; \mathbf{1}\!\left[\hat{f}^{\mathrm{LLM}}_i \text{ has no pathological outputs}\right].
  \label{eq:trust}
\end{equation}
Pathological outputs include: NaN values, infinite predictions, and division-by-zero
exceptions.  If the trust gate fails, the task is immediately routed to pure neural
fallback.

\paragraph{Stage 2 --- Transcendental Token Detection.}
Flag the formula if it contains any transcendental token from the set
$\{\exp, \log, \Phi, \sin, \cos, \mathrm{erf}\}$, indicating elevated structural
complexity.  Transcendental tokens are strongly predictive of stochastic instability
(Instability Index $\II > 0.05$) and trigger ensemble mode rather than pure LLM
routing.

\paragraph{Stage 3 --- Neural Degradation Probe.}
Compute the extrapolation degradation:
\begin{equation}
  \Delta_i = {R^{2,\mathrm{LLM}}_{\mathrm{train}}} - {R^{2,\mathrm{LLM}}_{\mathrm{edge}}}.
  \label{eq:degradation}
\end{equation}
A large $\Delta_i$ (above a threshold $\tau$, calibrated per domain category)
indicates that the LLM formula degrades rapidly near the training boundary,
suggesting it will fail further into the extrapolation region.  When $\Delta_i > \tau$,
ensemble mode is activated regardless of Stage 2 outcome.

\paragraph{Stage 4 --- Out-of-Distribution Detection.}
Estimate whether the test inputs are likely to be out of the training distribution:
\begin{equation}
  p_{\mathrm{OOD}} = \hat{\Pr}\!\left(X^{\mathrm{test}} \notin \mathrm{range}(X^{\mathrm{train}})\right),
  \label{eq:ood}
\end{equation}
computed using a kernel-density estimator fitted to the training projections $z^{(j)}$.
High $p_{\mathrm{OOD}}$ upweights the LLM formula (which generalises structurally) in
the ensemble.

\paragraph{Stage 5 --- Routing Decision and Ensemble Weights.}
The full routing decision is:
\begin{equation}
  \mathrm{decision}_i = \begin{cases}
    \text{Neural Fallback} & \text{if Trust}_i^{\mathrm{LLM}} = 0 \\
    \text{Pure LLM}        & \text{if } {R^{2,\mathrm{LLM}}_{\mathrm{train}}} \ge 0.95 \;\land\; \Delta_i < \tau \\
    \text{Ensemble}        & \text{if } {R^{2,\mathrm{LLM}}_{\mathrm{train}}} \ge 0.5 \\
    \text{Neural Fallback} & \text{otherwise}
  \end{cases}
  \label{eq:routing}
\end{equation}

When in ensemble mode, the distance-aware weight assigned to the LLM formula is:
\begin{equation}
  \wLLM = \alpha + (1 - \alpha)\cdot p_{\mathrm{OOD}},
  \label{eq:wllm}
\end{equation}
with $\wNN = 1 - \wLLM$ and $\alpha \in [0,1]$ a base weight (set to $\alpha=0.3$
throughout).  The final ensemble prediction is:
\begin{equation}
  \hat{y} = \wLLM \cdot \hat{f}^{\mathrm{LLM}}_i(x) + \wNN \cdot f^{\mathrm{NN}}_i(x').
  \label{eq:ensemble}
\end{equation}

\begin{remark}[On ensemble weights]
The weights in~\eqref{eq:wllm} are derived from $p_{\mathrm{OOD}}$, which is
estimated entirely from the training data.  This ensures that no test information
leaks into the prediction, addressing the data-leakage correction described in
Section~\ref{sec:corrections}.
\end{remark}

\subsection{Unified Formula Executor}\label{subsec:unified-formula-executor}
\label{sec:executor}

All symbolic formulas---whether produced by the LLM or specified as ground truth---are
evaluated through a single unified executor.  This executor:
\begin{itemize}
  \item Parses expressions into a normalised AST using SymPy.
  \item Compiles to NumPy vectorised operations for batch evaluation.
  \item Applies consistent type coercion (all inputs cast to \texttt{float64}).
  \item Catches and quarantines division-by-zero and overflow exceptions.
  \item Returns NaN for malformed expressions, triggering the trust-gate failure path.
\end{itemize}
Using a single evaluation backend eliminates the inconsistencies that affected
prior versions of the pipeline, where the same formula could produce different
$\Rsq$ values depending on the evaluation backend.

% =============================================================================
\section{HypatiaX Architecture}
\label{sec:architecture}

\subsection{Design Principles}\label{subsec:design-principles}
\label{sec:arch_principles}

\begin{enumerate}
\item \textbf{Symbolic methods for discovery:} Evolutionary search with domain
  constraints finds mathematically valid expressions \citep{cranmer2023pysr}.
\item \textbf{Multi-layer validation:} Dimensional analysis, domain constraints,
  numerical stability, and extrapolation testing.
\item \textbf{LLMs as interfaces:} Natural language interpretation and optional
  warm-starting \citep{kambhampati2024llms}.
\item \textbf{Failure transparency:} Invalid expressions are explicitly rejected.
\end{enumerate}

\subsection{Assumptions}\label{subsec:assumptions}
\label{sec:assumptions}

\begin{enumerate}
\item \textbf{Finite operator set:}
  $\mathcal{O} = \{+,-,\times,\div,\exp,\log,\sin,\cos,x^n,\sqrt{x}\}$.
\item \textbf{Correctly specified dimensional metadata.}
\item \textbf{Representative domain sampling.}
\item \textbf{Scale boundedness:} characteristic scales spanning
  $<6$ orders of magnitude.
  This threshold is derived from a single failure case (gravitational force
  constant $G = 6.674 \times 10^{-11}$) and should be treated as a heuristic
  rather than a validated boundary.
\item \textbf{Low measurement noise:}
  $\sigma = 0.05 \cdot \mathrm{std}(y)$ in all main experiments.
\end{enumerate}

\subsection{Five-Stage Routing Architecture Overview}\label{subsec:fivestage-routing-architecture-overview}
\label{sec:validation_framework}

Stage~1--2: data ingestion and preprocessing.
Stage~3 (optional): LLM proposes candidate expressions; on the v3.0 DeFi
benchmark, 68 of 74 tasks are routed to this path, delivering a $1.73\times$
median speedup over the neural baseline on those cases
(Section~\ref{sec:timing}).
Stage~4: PySR evolutionary search \citep{cranmer2023pysr}.
Stage~5: four-stage validation cascade (Section~\ref{sec:validation_detail})
producing the accepted expression with natural-language interpretation.

\subsection{System Variants}\label{subsec:system-variants}
\label{sec:arch_variants}

\textsc{HypatiaX} provides three architectural variants used as ablation
configurations:

\paragraph{Hybrid v50\_2 --- System 1 (extrapolation-focused).}
Integrates all five stages.  Achieved near-zero extrapolation error on the
Core~15 benchmark (Section~\ref{sec:five_systems}).

\paragraph{System 2 (pure symbolic with validation optimisation).}
Symbolic discovery without LLM guidance.

\paragraph{System 3 (LLM + symbolic fallback for robustness).}
Achieved $R^2 = 1.000$ on all 30 interpolation tests through automatic
fallback from low-confidence LLM outputs.

\subsection{Hybrid DeFi Variant: Architectural Specification}\label{subsec:hybrid-defi-variant-architectural-specif}
\label{sec:hybrid_defi_spec}

The Hybrid DeFi system is a domain-tuned variant of Hybrid~v50\_2 with the
following differences:

\begin{enumerate}
\item \textbf{Routing cascade (Fixes~0--5):} Six targeted routing improvements
  were developed on the DeFi development set (see Supplementary~A for full
  implementation details).  Fix~0 corrects data-generation bugs in zero-variance
  cases.  Fixes~1--4 replace in-distribution $R^2$ routing with uncertainty-aware
  ensemble voting:
  \[
    \hat{y} = w_\mathrm{LLM}\,\hat{y}_\mathrm{LLM}
             + w_\mathrm{SR}\,\hat{y}_\mathrm{SR}
             + w_\mathrm{NN}\,\hat{y}_\mathrm{NN},
    \quad w_k = \frac{\mathrm{conf}_k}{\sum_j \mathrm{conf}_j}.
  \]
  Fix~5 (unified formula evaluator) replaces the absolute SS threshold
  (\texttt{ss\_tot > 1e-10}) with a relative threshold, correcting the
  measurement bug described in Section~\ref{sec:r2_bugfix}.
\item \textbf{Domain operator set:} Operators augmented with financial
  primitives: $\max(0,\cdot)$, $\min(\cdot,\cdot)$, and
  $\ln(1+|\cdot|)$ (log-return normalisation).
\item \textbf{Confidence threshold:} Trust gate threshold $R^2_{\mathrm{train}} \geq 0.5$
  (v3.0 correction; raised from 0.1 in prior versions), with additional
  pathological-output detection.
\end{enumerate}

All other components (PySR configuration, four-layer validation, LLM model and
temperature) are identical to Hybrid~v50\_2.

\subsection{Symbolic Discovery Core}\label{subsec:symbolic-discovery-core}
\label{sec:arch_discovery}

We employ PySR \citep{cranmer2023pysr} for multi-objective evolutionary search:
\begin{align}
\text{Minimize:} \quad
& \mathcal{L}(E,\mathcal{D})
  = \frac{1}{n}\sum_{i=1}^n (E(x_i) - y_i)^2, \label{eq:loss}\\
\text{Minimize:} \quad
& \mathrm{Complexity}(E) = \sum_{t \in E} w(t), \label{eq:complexity}\\
\text{Subject to:} \quad
& E \in \mathcal{C}. \label{eq:constraints}
\end{align}
The algorithm maintains a Pareto front over prediction error and expression
complexity.

\subsection{Multi-Layer Validation Framework}\label{subsec:multilayer-validation-framework}
\label{sec:validation_detail}

The four validation layers target distinct failure modes:

\begin{itemize}
\item $V_1$ (12\% error detection): Dimensional consistency.
\item $V_2$ (18\%): Domain constraint verification.
\item $V_3$ (23\%): Numerical stability analysis.
\item $V_4$ (47\%): Predictive accuracy under extrapolation.
\end{itemize}

\begin{remark}[Validation-layer point estimates]
\label{rem:validation_ci}
The error detection rates above are point estimates over the Core~15 benchmark
($n=15$ equations).  Each rate carries an approximate $\pm$10--15\,pp
uncertainty (Wilson score interval, $n=15$).  These figures should be
interpreted as indicative of the relative contribution of each layer rather
than as precisely calibrated detection rates.
\end{remark}

\begin{finding}[Empirical Error Reduction via Cascaded Validation]
\label{prop:validation_reduction}
On the evaluated benchmark, the proportion of incorrect expressions surviving
all four layers is substantially lower than those surviving prediction-accuracy
validation alone.  No incorrect expression passed all four layers.
\end{finding}

% =============================================================================
\section{Experimental Setup}
\label{sec:setup}

\subsection{Benchmark Summary}\label{subsec:benchmark-summary}
\label{sec:setup_summary}

We evaluate HypatiaX v3.0 on the benchmark described in Section~\ref{sec:benchmark}:
74 DeFi tasks, all tractable (0 intractable cases), covering pricing, risk, and
liquidity modelling.  The fixed denominator for all aggregate statistics is 74;
NaN and failure outputs count as zero $\Rsq$ toward this denominator.

\subsection{Methods Compared}\label{subsec:methods-compared}
\label{sec:setup_methods}

Three methods are evaluated on each task:

\paragraph{Pure LLM (Symbolic Regression Baseline).}
The LLM generates a symbolic formula from the task description.  The formula is
executed directly with no numerical fitting or parameter optimisation.  All 30
independent stochastic runs use the same prompt; run-to-run variation arises
solely from temperature sampling.  \emph{For single-run evaluation (the standard
comparison with neural methods), one run is sampled uniformly at random.}

\paragraph{Neural Network (MLP Baseline).}
The residual MLP described in Section~\ref{sec:nn} is trained on augmented data
$(X, y) \cup (1.2X, 1.2y)$.  Training is subject to a 120-second wall-clock
limit per case.  Fixed random seeds are used throughout to ensure reproducibility.
The neural baseline does \emph{not} use LLM augmented features (Eq.~\ref{eq:augmented_input}),
ensuring a clean comparison of pure neural versus hybrid methods.

\paragraph{HypatiaX Hybrid (Proposed).}
The full five-stage routing and ensembling system described in Section~\ref{sec:routing}.
The hybrid method uses the same LLM formula and neural network as above, combined via
the distance-aware ensemble (Eq.~\ref{eq:ensemble}).  When neural fallback is triggered,
the neural network is retrained on the augmented input $(x')$ from scratch; this
retraining cost is fully included in the reported wall-clock time.

\subsection{Evaluation Metrics}\label{subsec:evaluation-metrics}
\label{sec:setup_metrics}

We report the following metrics for each method:
\begin{itemize}
  \item \textbf{Median test $\Rsq$}: robust measure of central tendency, less
    sensitive to catastrophic failures than the mean.
  \item \textbf{Mean test $\Rsq$ (clipped)}: arithmetic mean of $\min(\Rsq, 1)$,
    clipped at $-10$ to prevent extreme failures from dominating.
  \item \textbf{Near-perfect success rate} ($\Rsq > 0.99$): fraction of tasks where
    the method achieves near-exact formula recovery.
  \item \textbf{High-accuracy rate} ($\Rsq > 0.9$): fraction of tasks with strong
    predictive performance.
  \item \textbf{Catastrophic failure rate} ($\Rsq < -10$): fraction of tasks with
    prediction worse than the mean by a factor of more than 10.
  \item \textbf{Extrapolation gap} ($\Rtrain - \Rtest$): measures how much
    performance degrades from training to test.
  \item \textbf{Stability score} ($\Rtest / \Rtrain$): ratio of test to training
    performance; a value near 1.0 indicates consistent generalisation.
\end{itemize}

\paragraph{Fixed-denominator reporting.}
All percentages (e.g.\ $\Rsq > 0.99$ rate) use the fixed denominator of 74 across
all three methods.  A task that produces NaN or fails to execute contributes 0 to
the numerator of success rates.  This ensures that the three methods' rates are
directly comparable and that avoiding evaluation failures is not rewarded.

\subsection{Runtime Measurement}\label{subsec:runtime-measurement}
\label{sec:timing_setup}

Wall-clock time is recorded per task and per method.  For the hybrid system,
the reported time \emph{fully accounts} for all LLM inference time plus any
neural network retraining that occurs during fallback.  We do not separate
LLM and NN costs in the aggregate reported time; they are summed.

We separately report the speedup for the LLM-routed subset (cases where the
hybrid routes to pure LLM or ensemble, without triggering full NN retraining).
This subset comprises $n=68$ of 74 tasks and represents the cleanest comparison
for assessing the efficiency of symbolic routing.

\subsection{Implementation Details}\label{subsec:implementation-details}
\label{sec:setup_impl}

All experiments are implemented in Python~3.12 using:
\begin{itemize}
  \item \textbf{NumPy}: array operations and numerical evaluation.
  \item \textbf{PyTorch}: neural network training.
  \item \textbf{SciPy}: kernel-density estimation for OOD detection; PCA via
    \texttt{sklearn.decomposition.PCA}.
  \item \textbf{SymPy}: symbolic expression parsing and AST analysis.
\end{itemize}

Neural networks are trained with fixed random seeds (seed 42 for all cases) to
ensure full reproducibility of the NN baseline.  LLM formula generation uses a
deterministic prompting framework with retry logic (up to 3 retries) for API
stability, but temperature sampling is not disabled; this is the source of
run-to-run variability measured by the Instability Index.

% =============================================================================
\section{Results}
\label{sec:results}

\subsection{Five-System Comparative Analysis (Feynman Core-15)}\label{subsec:fivesystem-comparative-analysis-feynman-}
\label{sec:five_systems}

Table~\ref{tab:five_systems_full} compares the five system configurations on the
Core-15 Feynman benchmark at $2\times$ training range.

\begin{table}[htbp]
\centering
\begin{threeparttable}
\caption{Five-System Comparison: Extrapolation Error vs.\ Interpolation $R^2$.}
\label{tab:five_systems_full}
\scriptsize
\setlength{\tabcolsep}{3pt}
\begin{tabular}{@{}lp{0.6cm}p{1.3cm}p{1.3cm}p{1.4cm}p{0.7cm}p{1.4cm}@{}}
\toprule
\textbf{System} & \textbf{n} & \textbf{Extrap.} & \textbf{Extrap.} &
  \textbf{Train $R^{2}$} & \textbf{Std} & \textbf{Design} \\
 & & \textbf{Median (\%)} & \textbf{Mean (\%)} & \textbf{Mean} & & \textbf{Focus} \\
\midrule
\textbf{Hybrid v50\_2}\tnote{a} & 14 & \textbf{0.0} & \textbf{0.0} & 0.931\tnote{b} & --- & Extrapolation \\
Neural Network & 13 & 86.7 & 1231.0\tnote{c} & 0.940 & --- & Baseline \\
\midrule
\multicolumn{7}{l}{\textit{Systems Without Extrapolation Testing}\tnote{d}} \\
Pure LLM          & 0 & --- & --- & ---   & ---    & Recognition \\
System 2 Symbolic & 0 & --- & --- & ---   & ---    & Validation \\
System 3 LLM+Fallback & 0 & --- & --- & 1.000 & 0.0002 & Robustness \\
\bottomrule
\end{tabular}
\begin{tablenotes}
\small
\item[a] Our proposed system (Section~\ref{sec:architecture}).
\item[b] Mean $R^2 = 0.931$; median $R^2 = 1.000$ ($n=15$).
\item[c] Mean 1231\% at $2\times$ training range ($n=13$).
\item[d] Pure LLM: formula-to-callable compilation failures prevented
  extrapolation test execution.  Mann-Whitney $U=0$, $p=1.11\times10^{-6}$
  (Hybrid v50\_2 vs.\ NN).
\end{tablenotes}
\end{threeparttable}
\end{table}

\begin{result}[Five-System Extrapolation Hierarchy]
\label{thm:five_system_hierarchy}
Across 14 ground-truth equations at $2\times$ training range, Hybrid~v50\_2
achieves near-zero error (median $< 10^{-12}$ relative) while neural networks
exhibit mean error 1231\% (95\% CI: [1087\%, 1456\%], $n=13$), with no
distributional overlap: Mann-Whitney $U=0$, $p < 10^{-6}$.
\end{result}

\subsection{Overall Extrapolation Performance (DeFi v3.0)}\label{subsec:overall-extrapolation-performance-defi-v}
\label{sec:results_defi}

Table~\ref{tab:main_results} presents the primary results across all 74 benchmark tasks.

\begin{table}[h]
\centering
\caption{Aggregate extrapolation performance on the HypatiaX DeFi Benchmark (74 tasks).
All $\Rsq$ values clipped to $[-10,1]$; fixed denominator of 74.
Catastrophic: $\Rsq < -10$. \textbf{HypatiaX's $89.2\,\%$ figures are
uncorrected} for the hybrid decision-attribution bug
(\S\ref{sec:hybrid-attribution-bug}); the corrected overall near-perfect
rate is $60.8\,\%$ (below Pure LLM's $62.2\,\%$). The Catastrophic column
(0) is unaffected --- see the abstract's footnote for why.}
\label{tab:main_results}
\begin{tabular}{lrrrrr}
\toprule
Method & Median $\Rsq$ & Mean $\Rsq$ & $>0.99$ (\%) & $>0.9$ (\%) & Catastrophic \\
\midrule
Pure LLM      & 1.0000 & $-0.7571$ & 62.2 & 62.2 & 6 \\
Neural MLP    & $-0.4675$ & $-0.9482$ & 5.4  & 12.2 & 0 \\
HypatiaX      & 1.0000 & $+0.8721$ & \textbf{89.2} & \textbf{89.2} & \textbf{0} \\
\bottomrule
\end{tabular}
\end{table}

\begin{figure}[htbp]
\centering
\includegraphics[width=\linewidth]{fig18_r2_heatmap_improved}
\caption{Mean far-extrapolation ($5\times$) $\Rsq$ by formula domain and
difficulty tier (Easy/Med/Hard), Core-15 benchmark. Dashes indicate no
equation exists for that domain--difficulty combination in this benchmark.
HypatiaX is not uniformly better than PySR-only: it improves on
Chemistry (Easy: $0.90$ vs.\ $-12.56$) but produces $-\infty$ --- a genuine
evaluation crash, not merely a poor fit --- on DeFi AMM (Hard) and Physics
(Hard), where PySR-only instead returns a finite (if poor, $-78.04$, or
perfect, $1.00$) value. See \texttt{unconfirmed\_paper\_claims.tex} for the
per-equation breakdown underlying this aggregate.}
\label{fig:r2_heatmap_clipped}
\end{figure}

\begin{figure}[htbp]
\centering
\includegraphics[width=\linewidth]{fig09_r2_heatmap_regimes}
\caption{$\Rsq$ heatmap by equation (rows) and extrapolation regime
(columns: Train, Near $1.2\times$, Medium, Far $5\times$), Core-15
benchmark, on a shared colour scale saturating outside $[-1.5, 1.0]$.
\textbf{This figure, not Table~\ref{tab:llm_ablation}, is the source that
should be used to re-derive that withdrawn table} (see
\texttt{unconfirmed\_paper\_claims.tex}, Item~1): for six of fifteen
equations, HypatiaX Hybrid returns \texttt{nan} (Henderson-Hasselbalch,
Logistic Growth, Michaelis-Menten, Rate Law) or $-\infty$ (Constant Product,
Gravitational Force) beyond the Train column --- missing or crashed
evaluations, not $\Rsq$ scores of zero, and not to be confused with a poor
fit. PySR-only completes every cell; its most severe values are Impermanent
Loss (Far $=-157$) and Arrhenius (Far $=-13$). Notably, on Arrhenius,
HypatiaX Hybrid's completed Far value ($+0.90$) is a \emph{success}, the
opposite of the ``both collapse identically'' claim withdrawn from
\S\ref{sec:arrhenius_failure}. HypatiaX Hybrid's own most severe completed
value is Impermanent Loss (Far $=-8259$).}
\label{fig:r2_heatmap_raw}
\end{figure}

Several key observations emerge from Table~\ref{tab:main_results}.

\paragraph{Bimodal LLM behaviour.}
The pure LLM baseline achieves a median $\Rsq$ of $1.00$, which suggests strong
performance on a typical case.  However, the mean $\Rsq = -0.7571$ reveals a
stark bimodal distribution: the majority of tasks are recovered exactly, but
6 tasks suffer catastrophic failures ($\Rsq < -10$), dragging the mean far below
zero.  Moreover, the $>0.9$ success rate of 62.2\,\% (using a fixed denominator
of 74) reflects that a substantial fraction of tasks fall outside perfect recovery
even when failures are not catastrophic.

\paragraph{Neural network extrapolation collapse.}
The neural MLP baseline performs poorly under the aggressive extrapolation split,
with a median $\Rsq$ of $-0.47$---systematically \emph{worse} than predicting
the mean.  This is consistent with the well-established failure of MLPs under
distribution shift.  Despite data augmentation ($1.2\times$ scaling), the network
does not recover the functional form of the underlying formula and instead learns
a locally accurate but extrapolation-sensitive approximation.  Notably, the neural
baseline has zero catastrophic failures: it always produces finite predictions,
just highly inaccurate ones.

\paragraph{HypatiaX eliminates catastrophic failures.}
The hybrid method achieves a median $\Rsq$ of $1.00$ (matching the best LLM cases)
while raising the near-perfect success rate from 62.2\,\% to \textbf{89.2\,\%}
\textbf{[flagged, see \S\ref{sec:hybrid-attribution-bug}: corrected 60.8\,\%]}
and eliminating all 6 catastrophic failures.  The mean $\Rsq$ improves from
$-0.7571$ (LLM) and $-0.9482$ (NN) to $+0.8721$---a reversal of sign that reflects
the elimination of catastrophic failures. (The elimination of catastrophic
failures itself is unaffected by the bug described in
\S\ref{sec:hybrid-attribution-bug}: the fabricated hybrid $\Rsq$ values are all
$\approx 1.0$, not catastrophic, so zero catastrophic failures remains correct.)

\paragraph{Statistical significance (Core-15 benchmark).}
Across the Core-15 benchmark equations, \textsc{HypatiaX} does not achieve
statistically significant higher far-extrapolation $\Rsq$ than \textsc{PySR}-only
(Mann-Whitney $U=126$, $p=0.2948$, $n=15$; rank-biserial $r=-0.12$).  The result
is stable across two hardware runs (Run~B: $U=101.5$, $p=0.6833$).
Domain-selective improvement is observed in Physics and DeFi Risk;
\textsc{PySR}-only is superior in Chemistry and DeFi AMM.
These null results on the Core-15 subset are expected given the full-benchmark
picture: significant gains emerge in the Feynman success-subset
(Mann-Whitney $U=81$, $p=0.00020$) and are concentrated on the hard transcendental
and DeFi instability cases discussed in Sections~\ref{sec:instability_results}
and~\ref{sec:feynman30}.

\subsection{Performance by Difficulty Level}\label{subsec:performance-by-difficulty-level}
\label{sec:difficulty}

Table~\ref{tab:difficulty} breaks down the near-perfect success rate ($\Rsq > 0.99$)
by difficulty tier.

\begin{table}[h]
\centering
\caption{Near-perfect success rate ($\Rsq > 0.99$) by difficulty.
Fixed denominator per tier; LLM and Hybrid use single-run evaluation.
\textbf{The ``HypatiaX (\%)'' column is flagged}: it is computed from the
hybrid system's own \texttt{success}/$\Rsq$ fields without cross-checking the
\texttt{decision} field against the actual success of the sub-method it
names. The ``HypatiaX, corrected (\%)'' column instead treats a task as a
hybrid success only when the sub-method the \texttt{decision} field names
itself reports \texttt{success=True}. See \S\ref{sec:hybrid-attribution-bug}.}
\label{tab:difficulty}
\begin{tabular}{lrrrrr}
\toprule
Difficulty & $n$ & Pure LLM (\%) & HypatiaX (\%) & HypatiaX, corrected (\%) & Gain (pp), corrected \\
\midrule
Easy       & 24  & 87.5          & \textbf{100.0} & 87.5 & $0.0$ \\
Medium     & 29  & 58.6          & \textbf{93.1}  & 58.6 & $0.0$ \\
Hard       & 21  & 38.1          & \textbf{76.2}  & 33.3 & $-4.8$ \\
\midrule
Overall    & 74  & 62.2          & \textbf{90.5}  & 60.8 & $-1.4$ \\
\bottomrule
\end{tabular}
\end{table}

The gain is monotonically increasing with difficulty in the uncorrected figures:
$+12.5$ percentage points (pp) on easy tasks, $+31.1$\,pp on medium, and $+38.1$\,pp
on hard \textbf{[flagged: corrected gains are $0.0$/$0.0$/$-4.8$\,pp respectively
--- monotonicity does not survive correction, and the hard-tier figure becomes a
loss, not the largest gain; see \S\ref{sec:hybrid-attribution-bug}]}. Before
correction, this was read as evidence that
the hybrid routing mechanism is most effective precisely where the pure LLM is most
unreliable---on transcendental and multi-asset risk formulas; that reading does
not hold once decision-attribution is corrected for.

\paragraph{Easy tasks.}
On easy (algebraic) tasks, the LLM already achieves 87.5\,\% near-perfect recovery,
reflecting strong LLM capability on linear and rational formulas.  HypatiaX's
raw figure is 100\,\%, but this tier is not exempt from the decision-attribution
bug: 3 of the 24 easy-tier successes are fabricated (the named sub-method
itself failed), giving a corrected rate of 87.5\,\%
\textbf{[corrected, \S\ref{sec:hybrid-attribution-bug}]} --- exactly matching
the pure-LLM baseline, i.e.\ once fabricated cases are removed, HypatiaX shows
\emph{no measured gain} on easy tasks in this run.

\paragraph{Medium tasks.}
On medium tasks, the LLM success rate drops to 58.6\,\%, reflecting the additional
complexity of nonlinear and exponential formulas.  HypatiaX's raw figure is 93.1\,\%
\textbf{[flagged: corrected 58.6\,\%, \S\ref{sec:hybrid-attribution-bug}]}
--- 10 of the 27 raw medium-tier successes are fabricated, and the corrected
rate again matches the pure-LLM baseline exactly, indicating no measured gain
on this tier once the bug is corrected for.

\paragraph{Hard tasks.}
On hard (transcendental and multi-asset) tasks, the pure LLM achieves only 38.1\,\%
near-perfect recovery---barely above chance for a binary success criterion.  This
reflects the stochastic collapse documented in the Instability Index analysis:
Black--Scholes, Theta, and correlated portfolio risk formulas are recovered correctly
on some runs but fail completely on others.  HypatiaX's raw figure is 76.2\,\%
\textbf{[flagged: corrected 33.3\,\%, \S\ref{sec:hybrid-attribution-bug}]} ---
9 of the 16 raw hard-tier successes are fabricated. Unlike the raw figures,
the corrected rate on hard tasks (33.3\,\%) is now \emph{below} the pure-LLM
baseline (38.1\,\%), a $-4.8$\,pp gap: once fabricated successes are removed,
HypatiaX does not outperform the pure LLM on hard tasks in this run.
The fabricated cases are spread across all three tiers (3 easy, 10 medium, 9
hard; see \S\ref{sec:hybrid-attribution-bug} for the full per-task list) ---
this is \emph{not} concentrated in the hard tier as an earlier draft of this
analysis stated, and that concentration claim is retracted.

\subsection{The Hybrid Decision-Attribution Bug: Quantification}
\label{sec:hybrid-attribution-bug}

\paragraph{The bug.} Every hybrid-system record in the v3.0 result JSON carries
three fields relevant to success accounting: \texttt{hybrid.success},
\texttt{hybrid.test\_r2}, and \texttt{hybrid.decision} (one of \texttt{llm},
\texttt{nn}, or \texttt{nn\_fallback}, naming which sub-method's output the
hybrid system is reporting). All success-rate figures in this paper prior to
this section were computed from \texttt{hybrid.success}/\texttt{hybrid.test\_r2}
alone. This is incorrect whenever \texttt{hybrid.decision} names a sub-method
that itself failed: the pipeline does not appear to verify, at the point it
records \texttt{hybrid.success = True}, that the named sub-method actually
produced a valid fit. Concretely, we observe records with
\texttt{hybrid.decision = "llm"}, \texttt{hybrid.test\_r2} $\approx 1.0$, and
\texttt{hybrid.success = True}, for tasks where \texttt{pure\_llm.success =
False} (the LLM method failed to execute and returned \texttt{NaN}) on that
same task in the same run.

\paragraph{Methodology.} We define a \emph{corrected} hybrid success indicator:
a task counts as a hybrid near-perfect success only if (a)
\texttt{hybrid.test\_r2} $> 0.99$ \emph{and} (b) the sub-method named by
\texttt{hybrid.decision} itself reports \texttt{success = True} on that task
(\texttt{pure\_llm.success} when \texttt{decision = "llm"};
\texttt{neural\_network.success} when \texttt{decision} is \texttt{nn} or
\texttt{nn\_fallback}). This is a strictly conservative correction: it does
not penalise any task where the hybrid system's own accounting is internally
consistent, and only removes successes that are attributable to a sub-method
the pipeline itself marked as failed.

\paragraph{Magnitude (seed 42, all 74 tasks).} Applying this correction to
\texttt{hypatiax\_defi\_benchmark\_v3\_results\_seed42.json} yields 22 fabricated
near-perfect successes out of the 67 the pipeline reports. All 22 route through
\texttt{decision = "llm"} --- none occur on the \texttt{nn}/\texttt{nn\_fallback}
path in this run. The corrected overall near-perfect rate is $45/74 = 60.8\,\%$,
versus the $67/74 = 90.5\,\%$ the uncorrected pipeline reports for this run.

\paragraph{Reconciling with the previously published $89.2\,\%$/$66$-of-$74$
figure.} This is one task fewer than the $67/74$ recomputed here from
\texttt{hypatiax\_defi\_benchmark\_v3\_results.json} (seed~42). We checked
whether this is a threshold/rounding artefact by examining the seven tasks
in this seed's run that do \emph{not} clear $\Rsq > 0.99$: their
\texttt{hybrid.test\_r2} values are $0.771, 0.542, 0.659, 0.343, 0.0,
-0.775, -0.898$ --- the closest is $0.22$ below the cutoff, none are
borderline, and none would flip under any reasonable rounding convention.
This rules out rounding as the explanation. The remaining, better-supported
explanation is ordinary cross-seed nondeterminism in the pipeline itself
(the LLM proposal step and PySR's evolutionary search are both stochastic):
a run on a different seed than 42 --- plausibly whichever seed produced the
originally published Table~\ref{tab:difficulty} --- could easily have one
fewer task clear the threshold than seed~42 does, without any task's value
being close to the boundary in \emph{either} run. \textbf{What this closes:}
the one-task gap is not a computational or reconciliation error in this
correction, and does not need a rounding-based explanation. \textbf{What
remains open:} we have not identified which specific seed (of $\{42, 99,
123, 777, 2024\}$) produced the exact JSON underlying the originally
published $66/74$ figure, so that figure's per-task composition is not
independently confirmed here. This does not affect the conclusion either
way: every seed we have examined shows the same decision-attribution bug,
and the corrected $60.8\,\%$ figure is computed independently of which seed
is treated as canonical. Table~\ref{tab:hybrid-bug-breakdown}
gives the full breakdown.

\paragraph{Independent cross-seed confirmation on a second, smaller sample.}
A separate 5-task Portfolio-focused sample
(\texttt{hypatiax\_defi\_benchmark\_v3\_results\_seed\{42,99,123,777,2024\}.json})
lets us check whether at least one fabricated-success case is deterministic
across seeds rather than a single-run artefact. \emph{Portfolio Expected
Shortfall for correlated assets} shows \texttt{hybrid.decision = "llm"},
\texttt{hybrid.success = True}, \texttt{hybrid.test\_r2 = 1.0} in
\textbf{all five seeds}, while \texttt{pure\_llm.success = False} in all
five as well --- the identical fabricated-success pattern, reproducing
deterministically 5/5, ruling out seed-to-seed noise as an explanation for
\emph{this particular} case. The other four tasks in this sample show no
such inconsistency in any seed. This sample is too small (5 tasks, a
different subset from the 74-task suite above) to speak to the 66-vs-67
reconciliation in the previous paragraph, but it does independently confirm
that the decision-attribution bug itself is a deterministic property of
specific tasks, not merely a stochastic artefact of one run.

\begin{table}[h]
\centering
\caption{Hybrid decision-attribution bug: fabricated near-perfect successes by
difficulty tier, computed from \texttt{hypatiax\_defi\_benchmark\_v3\_results\_seed42.json}.}
\label{tab:hybrid-bug-breakdown}
\begin{tabular}{lrrrr}
\toprule
Difficulty & $n$ & Reported HypatiaX successes & Fabricated & Corrected successes \\
\midrule
Easy    & 24 & 24 & 3 & 21 \\
Medium  & 29 & 27 & 10 & 17 \\
Hard    & 21 & 16 & 9 & 7 \\
\midrule
Overall & 74 & 67 & 22 & 45 \\
\bottomrule
\end{tabular}
\end{table}

\paragraph{Fabricated-success task list (seed 42, 22 tasks, all
\texttt{decision = "llm"}).} \emph{Easy (3):} Loan-to-Value, LP share
percentage, Slashing penalty. \emph{Medium (10):} Liquidation Price Long,
Liquidation Price Short, Effective Leverage, Capital efficiency, Price
slippage percentage, Utilization rate of DeFi, Impermanent loss percentage,
Collateral ratio, Staking reward for fixed lock-up, Concentrated liquidity
position width. \emph{Hard (9):} Component ES, Liquidation price for
leveraged long, Liquidation price for leveraged short, Maximum safe
leverage, Required collateral, Portfolio Expected Shortfall for correlated,
Call option intrinsic, Simple options moneyness, Uniswap V3 virtual.
The last of these is confirmed independently across all five available
seeds (42, 99, 123, 777, 2024): in every seed, \texttt{pure\_llm.success =
False} (\texttt{NaN} train/test $\Rsq$) while \texttt{hybrid.decision = "llm"}
and \texttt{hybrid.test\_r2 = 1.0}, a deterministic pattern that rules out
run-to-run noise as the explanation.

\paragraph{What this does and does not affect.} The corrected figures leave
one of this paper's central qualitative claims intact: HypatiaX's own
reported output still shows no catastrophic failures ($\Rsq < -10$), since
every fabricated success is reported near $\Rsq \approx 1.0$, not a
catastrophic value --- the underlying catastrophic failures (in the masked
sub-method, typically \texttt{pure\_llm}) are hidden by the bug, not absent
from the computation; see the abstract's footnote on this distinction.

The claim that HypatiaX outperforms the pure LLM baseline at every
difficulty tier \textbf{does not survive correction} and is retracted. At
the corrected rate, HypatiaX ties the pure-LLM baseline on Easy ($87.5\,\%$
vs.\ $87.5\,\%$) and Medium ($58.6\,\%$ vs.\ $58.6\,\%$), and underperforms
it on Hard ($33.3\,\%$ vs.\ $38.1\,\%$, a $-4.8$\,pp gap). Overall, the
corrected HypatiaX rate ($60.8\,\%$) is \emph{below} the pure-LLM baseline
($62.2\,\%$), a $-1.4$\,pp gap. This reverses the paper's original headline
comparison (previously reported as $+27.0$\,pp overall, $+38.1$\,pp on
Hard) rather than merely shrinking it. We do not know, from the evidence in
this section alone, whether this reflects a genuine absence of routing
benefit once decision-attribution is correctly accounted for, or whether
the hybrid routing mechanism itself needs to be re-evaluated on tasks where
its named sub-method actually succeeds (i.e.\ whether HypatiaX would show a
real gain if the router made better decisions rather than merely reporting
them incorrectly). This distinction matters for interpretation and is not
resolved here. This is a materially different picture from
what Section~\ref{subsec:performance-by-difficulty-level} originally
reported, and downstream claims that cite a HypatiaX-over-LLM gain on any
tier should be read against the corrected figures above, not the
originally reported ones, until the underlying pipeline bug is fixed and
the benchmark is rerun.

\subsection{Runtime Analysis}\label{subsec:runtime-analysis}
\label{sec:timing}

Table~\ref{tab:timing} presents the full wall-clock timing data from the v3.0
benchmark run.

\begin{table}[h]
\centering
\caption{Wall-clock time per task (seconds). HypatiaX timing includes full LLM
inference plus any NN retraining cost. Speedups computed relative to Neural MLP.}
\label{tab:timing}
\begin{tabular}{lrrrr}
\toprule
Method & Mean (s) & Median (s) & $n$ & vs.\ NN (mean) \\
\midrule
Pure LLM   & 11.4 & 10.3 & 74 & $3.80\times$ slower \\
Neural MLP &  3.0 &  2.7 & 74 & --- (baseline) \\
HypatiaX   &  6.8 &  1.7 & 74 & $2.30\times$ slower (mean) \\
            &      &      &    & $1.64\times$ faster (median) \\
\midrule
HypatiaX (LLM-routed only) & --- & --- & 68 & $1.73\times$ faster \\
\bottomrule
\end{tabular}
\end{table}

\paragraph{Interpretation.}
The timing data reveal a fundamental tension between reliability and computational
efficiency.

The hybrid method has a \textbf{median time of 1.7\,s}---substantially lower than
both the pure LLM (10.3\,s) and the neural network (2.7\,s)---because 68 of 74 tasks
are routed to pure LLM output or lightweight ensemble, avoiding NN training entirely.
For these 68 LLM-routed cases, the median speedup over the neural network is
$\mathbf{1.73\times}$.

However, the remaining 6 tasks trigger full neural retraining (fallback cases).
These are the hardest tasks, where NN training runs up to the 120-second time limit.
The presence of these fallback cases causes the \textbf{mean time to rise to 6.8\,s},
which is $2.30\times$ \emph{slower} than the neural baseline mean of 3.0\,s.

\paragraph{Correcting the 73\,\% reduction claim.}
\label{sec:correction}

An earlier draft of this paper reported a ``73\,\% runtime reduction'' over the
neural baseline, corresponding to a $3.7\times$ speedup.  This claim
is \textbf{not supported} by the v3.0 benchmark data.  The measured speedups are:
\begin{itemize}
  \item \emph{Mean-based}: HypatiaX is $2.30\times$ \textbf{slower} than NN
    ($6.8$\,s vs.\ $3.0$\,s); this corresponds to a $-129.9$\,\% ``reduction'' (i.e.\ an increase).
  \item \emph{Median-based}: HypatiaX is $1.64\times$ faster than NN
    ($1.7$\,s vs.\ $2.7$\,s); this corresponds to a $39.0$\,\% reduction.
  \item \emph{LLM-routed subset} ($n=68$): $1.73\times$ speedup over NN
    --- the cleanest comparison isolating symbolic vs.\ learned computation.
\end{itemize}

The correct claim for this paper is therefore: \textbf{when the LLM formula is
trusted (68 of 74 cases, 92\,\%), HypatiaX achieves a $1.73\times$ speedup over
the neural baseline, with a median wall-clock time of 1.7\,s per task.}  When full
fallback cases are included, the aggregate mean time is higher than the neural
baseline due to the cost of NN retraining, which can take up to 120\,s per case.

\paragraph{Interpretation for deployment.}
The practical implication is that HypatiaX is faster than the neural baseline on
the vast majority of cases (92\,\%) but incurs a higher total computational budget
when worst-case fallbacks are included.  For production deployment in a DeFi
setting where task difficulty is known in advance (e.g.\ algebraic AMM formulas
vs.\ Black--Scholes), the routing decision can be pre-computed, eliminating fallback
overhead entirely and achieving the $1.73\times$ speedup universally.

\subsection{Portfolio Variance: Seed-Sweep Robustness Analysis}\label{subsec:portfolio-variance-seedsweep-robustness-}
\label{sec:portfolio_seed_sweep}

Symbolic regression methods based on evolutionary search are inherently stochastic.
To assess robustness, we evaluate both PySR-only and \textsc{HypatiaX} on the
Portfolio Variance task across five independent random seeds
$\{42, 99, 123, 777, 2024\}$.

\begin{table}[h]
\centering
\caption{Portfolio Variance seed-sweep results.
  \textbf{H recovers?}: exact closed-form formula recovered.
  \textbf{H wins?}: HypatiaX far-$\Rsq$ strictly greater than PySR-only.}
\label{tab:portfolio_seed_sweep}
\small
\begin{tabular}{lccccl}
\toprule
\textbf{Seed} & \textbf{P far-$\Rsq$} & \textbf{H far-$\Rsq$} &
\textbf{H formula} & \textbf{H recovers?} & \textbf{H wins?} \\
\midrule
42 (default) & $-21.004$ & $-0.023$  & linear    & No  & Yes \\
99           & $-1.226$  & $-15.191$ & linear    & No  & No  \\
123          & $-18.651$ & $-18.090$ & exp denom & No  & Yes \\
777          & $-0.438$  & $+1.000$  & exact     & Yes & Yes \\
2024         & $-12.109$ & $+1.000$  & exact     & Yes & Yes \\
\midrule
Mean         & $-10.686$ & $-6.261$  &           & \multicolumn{2}{l}{H: 4/5 wins, 2/5 exact} \\
\bottomrule
\end{tabular}
\end{table}

\begin{figure}[htbp]
\centering
\includegraphics[width=\linewidth]{fig1_seed_sweep}
\caption{Portfolio Variance seed sweep: Near, Medium, and Far $\Rsq$ for PySR-only
and HypatiaX across five random seeds. HypatiaX recovers exact
formulae on seeds~777 and~2024 (Far $\Rsq = 1.00$). On seeds~42 and~123, both
methods perform poorly at far extrapolation, though HypatiaX is somewhat
higher (seed~42: $\approx 0$ vs.\ $\approx -21$; seed~123: $\approx -18$ vs.\
$\approx -18.5$, a marginal gap). PySR-only's far-$\Rsq$ ranges from
$\approx -21$ (seed~42) to $\approx -0.5$ (seed~777) --- not uniformly
``large negative'' across seeds. Seed~99 is the single case where HypatiaX
underperforms PySR-only at far extrapolation ($\approx -15$ vs.\ $\approx
-1.2$).}
\label{fig:portfolio_seed_sweep}
\end{figure}

PySR-only fails to extrapolate correctly under \emph{all} tested seeds ($0/5$),
whereas \textsc{HypatiaX} recovers the exact closed-form portfolio-variance formula
in $40\%$ of runs ($2/5$, seeds 777 and 2024) and achieves strictly higher far-$\Rsq$
in 4 of 5 seeds.  A Bayesian superiority analysis yields
$P(\text{HypatiaX} > \text{PySR}) \approx 0.76$, indicating that \textsc{HypatiaX}
occupies a fundamentally different regime of symbolic discovery.

\textbf{Corrected attribution.}\footnote{An earlier draft erroneously attributed
far-$\Rsq=-18.1$ to the default seed~(42).  The verified figure is: H seed=42
far-$\Rsq=-0.023$; the $-18.1$ value belongs to seed~123.  Figures of $-882.9$ and
``seed=99\,$\to$\,$\Rsq=1.000$'' derive from a separate ablation experiment (Exp-1)
and do not appear in the seed-sweep JSON; they must not be cited as DeFi seed-sweep
results.}

% ── RF-01 Option B: PySR-only vs HypatiaX ablation (Core-15) ────────────────
\subsection{Ablation: PySR-only vs.\ HypatiaX (Core-15)}\label{subsec:ablation-pysronly-vs-hypatiax-core15}
\label{sec:ablation_core15}

Table~\ref{tab:llm_ablation} presents the full per-equation comparison between
PySR-only ($P$) and HypatiaX ($H$) on the Core-15 benchmark across three
extrapolation ranges (near $1.2\times$, canonical, far $5\times$).  The
Mann-Whitney U-test on far-$\Rsq$ vectors ($n=15$): $U=126$, $p=0.2948$
(two-sided); rank-biserial $r=-0.12$.  Run~B: $U=101.5$, $p=0.6833$.
The null result is stable across two hardware runs, indicating that the
HypatiaX advantage is domain-selective rather than universal.  Specifically:
\textbf{Physics} (Kinetic Energy, Gravitational Force, Ideal Gas Law) and
\textbf{DeFi Risk} (Value at Risk, Liquidation Price, Portfolio Variance) favour
HypatiaX; \textbf{Chemistry} (Arrhenius, Henderson-Hasselbalch) and
\textbf{DeFi AMM} (Impermanent Loss) favour PySR-only.  The null aggregate
$p$-value reflects this domain cancellation: gains on Physics/Risk are offset by
losses on Chemistry/AMM.  We report this honestly; the DeFi benchmark
($n=74$, §\ref{sec:results}) shows the effect size when evaluated on
HypatiaX's target domain.

\textbf{This table has been rebuilt from the underlying \texttt{\_merged.json}
data} (via \texttt{table6\_vs\_real\_data\_comparison.xlsx}), replacing the
withdrawn version. Six equations show \textit{nan} or $-\infty$ for HypatiaX
beyond the Train column --- these are missing or crashed evaluations, not
$\Rsq$ scores, and are reported as such rather than omitted. See
\texttt{unconfirmed\_paper\_claims.tex}, Item~1, for the full discrepancy
history against the originally-published (withdrawn) numbers.

\begin{table*}[htbp]
\centering
\caption{LLM Ablation: PySR Alone (P) vs.\ HypatiaX (H) on Core~15 ---
\textbf{rebuilt from \texttt{\_merged.json}} (real data, not the withdrawn
manuscript transcription). Extrap columns show $R^2$ at near ($1.2\times$),
medium (canonical), and far ($5\times$) out-of-distribution ranges.
\textit{nan}: evaluation never completed. $-\infty$: evaluation crashed.}
\label{tab:llm_ablation}
\small\setlength{\tabcolsep}{3pt}
\resizebox{\linewidth}{!}{%
\begin{tabular}{llcccccccc}
\toprule
\textbf{Equation} & \textbf{Domain} & \multicolumn{2}{c}{\textbf{Train $R^2$}} & \multicolumn{2}{c}{\textbf{Near $R^2$}} & \multicolumn{2}{c}{\textbf{Med $R^2$}} & \multicolumn{2}{c}{\textbf{Far $R^2$}} \\
\cmidrule(lr){3-4}\cmidrule(lr){5-6}\cmidrule(lr){7-8}\cmidrule(lr){9-10}
 & & P & H & P & H & P & H & P & H \\
\midrule
Allometric Scaling & Biology & 0.9977 & 0.9979 & 1.0000 & 0.9824 & 1.0000 & 0.9502 & 1.0000 & $-0.0644$ \\
Arrhenius & Chemistry & 0.9974 & 0.9979 & 0.9779 & 0.9977 & 0.4636 & 0.9957 & $-12.5552$ & \textbf{0.9028} \\
Constant Product & DeFi AMM & 0.9982 & 0.9985 & 0.9996 & $-1.807\times10^{4}$ & 1.0000 & $-5.416\times10^{34}$ & 0.9996 & $-1.852\times10^{46}$ \\
Gravitational Force & Physics & 0.9976 & 0.9889 & 0.9973 & $-\infty$ & 0.9998 & $-\infty$ & 0.9973 & $-\infty$ \\
Henderson-Hasselbalch & Chemistry & 0.9339 & 0.9986 & 0.8004 & \textit{nan} & 0.9850 & \textit{nan} & 0.8004 & \textit{nan} \\
Ideal Gas Law & Physics & 0.9976 & 0.9980 & 1.0000 & 0.9988 & 1.0000 & 0.9999 & 1.0000 & 0.9993 \\
Impermanent Loss & DeFi AMM & 0.9971 & 0.9977 & 0.7703 & 0.9639 & $-2.7832$ & $-80.4327$ & $-157.0776$ & $-8259.3$ \\
Kinetic Energy & Physics & 0.9968 & 0.9968 & 1.0000 & 1.0000 & 1.0000 & 1.0000 & 1.0000 & 1.0000 \\
Liquidation Price & DeFi Risk & 0.9974 & 0.9977 & 1.0000 & 0.9998 & 1.0000 & 0.9999 & 1.0000 & 0.9993 \\
Logistic Growth & Biology & 0.9975 & 0.9976 & 0.9997 & \textit{nan} & 0.9998 & \textit{nan} & 0.9996 & \textit{nan} \\
Michaelis-Menten & Biology & 0.9975 & 0.9979 & 0.9890 & \textit{nan} & 0.9951 & \textit{nan} & \textbf{0.6471} & \textit{nan} \\
Portfolio Std Dev & DeFi Risk & 0.9975 & 0.9975 & 1.0000 & 1.0000 & 1.0000 & 1.0000 & 1.0000 & 1.0000 \\
Price Impact & DeFi AMM & 0.9976 & 0.9976 & 1.0000 & 1.0000 & 1.0000 & 1.0000 & 1.0000 & 1.0000 \\
Rate Law & Chemistry & 0.9977 & 0.9981 & 1.0000 & \textit{nan} & 1.0000 & \textit{nan} & 1.0000 & \textit{nan} \\
Value at Risk & DeFi Risk & 0.9979 & 0.9979 & 0.9999 & 1.0000 & 1.0000 & 1.0000 & 0.9999 & 1.0000 \\
\bottomrule
\end{tabular}%
}% end resizebox
\par\vspace{4pt}
{\footnotesize\noindent
P = PySR-only; H = HypatiaX (PySR + LLM warm-start). \textbf{Bold}: cases
where the real data reverses the withdrawn table's headline claim (Arrhenius
and Michaelis-Menten far-$R^2$ for the relevant column). H is \textit{nan}
or $-\infty$ (not evaluated, or crashed) on 6 of 15 equations beyond Train ---
Constant Product, Gravitational Force, Henderson-Hasselbalch, Logistic
Growth, Michaelis-Menten, Rate Law. Aggregate statistics (mean, Mann-Whitney)
are intentionally not recomputed here pending a decision on how to treat
\textit{nan}/$-\infty$ cells in a summary statistic; see
\texttt{unconfirmed\_paper\_claims.tex}, Item~1.
}
\end{table*}

\subsection{Feynman Extrapolation Benchmark ($n=30$)}\label{subsec:feynman-extrapolation-benchmark-n30}
\label{sec:feynman30}

To assess generalisation beyond DeFi, we evaluate \textsc{HypatiaX}
(\texttt{HybridDiscoverySystem\_v50\_2}) on 30 Feynman equations spanning
mechanics, thermodynamics, optics, electromagnetism, electrostatics,
magnetism, quantum mechanics, biology, chemistry, electrochemistry, and
probability.

\textbf{Split-protocol disclosure (FIX-C3).}\quad
Prior runs of the Feynman benchmark used a random 80/20 train/test split
(\texttt{sklearn.model\_selection.train\_test\_split}, \texttt{test\_size=0.2}),
whereas the DeFi benchmark uses a PCA-directed 40\,\%/60\,\% aggressive
extrapolation split throughout (Section~\ref{sec:split}).
The previously reported primary result (9/30 successes, Kaggle 4-vCPU run,
attributed to \texttt{fixc3\_baseline.json}) does not reproduce and is
withdrawn: that file contains no computed values
($\mathtt{n\_pass}{=}0$, $\mathtt{n\_total}{=}0$, $\mathtt{solve\_rate}{=}\mathtt{null}$)
and is an unmodified setup-script template stub, found byte-identical across
more than a dozen unrelated experiment directories rather than a value
computed for this experiment.

\textbf{Run provenance.}\quad
At least four distinct PCA-split run dates exist in this project's history;
Table~\ref{tab:provenance} discloses all four so that a reader encountering
any of the other three elsewhere cannot mistake them for the figures reported
in this section. Only the \textbf{July~22 / July~23, 2026} run pair is
canonical for Section~\ref{sec:feynman30}: July~22 is the random 80/20 split
source and July~23 is the PCA-directed 40/60 split source.

\begin{table}[h]
\centering
\caption{Every PCA-split run date identified in this project's history, and
  its status with respect to Section~\ref{sec:feynman30}. Only the
  July~22/23 pair is the source of the figures reported below.}
\label{tab:provenance}
\small
\begin{tabular}{lp{9.3cm}}
\toprule
Run date & Status \\
\midrule
June~4, 2026 & \textbf{Discarded, pre-fix.} \texttt{HybridDiscoverySystem}
scores 0/30, with all 30 equations flagged as environment/import failures
--- the method did not execute at all. Not used in this section. \\
June~29, 2026 & \textbf{Not located.} Named elsewhere as a canonical batch
but not recovered from any file set examined; its absence does not affect
the figures below, which are independently sourced from July~22/23. \\
June~30, 2026 & \textbf{Unverified, not used.} Source of a pooled
$n_{\text{pass}}{=}142$, $n_{\text{total}}{=}180$ figure that is, by its own
metadata, paired with the already-withdrawn ``9/30'' claim above and shares
that claim's known failure mode (arithmetic pooling of one run's records as
if two). Treated as unconfirmed pending a direct recount; not used in this
section. \\
\textbf{July~22 / July~23, 2026} & \textbf{Canonical.} Source of every
figure reported in this section. \\
\bottomrule
\end{tabular}
\end{table}

\paragraph{Threshold policy (stated once, applied everywhere).} Every
headline pass/fail figure in this paper --- abstract, main text, tables,
and conclusion alike --- uses one criterion: $\Rsq \geq 0.999999$
(``strict''), applied identically to both the random-split and PCA-split
runs. A result with $\Rsq \in [0.999, 0.999999)$ is disclosed separately
as a \emph{near-miss} and is never counted toward a headline pass figure
anywhere in this paper. Where a near-miss-inclusive (``loose'') tally is
reported below for transparency, it is always labeled explicitly by split
and convention (e.g.\ ``random-split, loose: 13/30'') and is never written
as a bare, unqualified fraction --- because, as it happens, the number 13
recurs in this section with three unrelated meanings (PCA-split strict
passes, random-split loose-convention passes, and random-split's genuine-fail
count) that must not be conflated. Any bare ``12/30'' or ``13/30'' elsewhere
in this paper refers to the strict headline figures defined here: random-split
strict $=12/30$, PCA-split strict $=13/30$.

Recomputed directly from raw per-test JSON records for the canonical
July~22/23 pair (not from a pre-merged summary file), and gated behind a
self-verification step --- exact reproduction of every already-scored
equation from the regenerated data, before any corrected null score is
trusted --- \textsc{HypatiaX} achieves \textbf{12/30 (40.0\,\%)} successes
($\Rsq \geq 0.999999$) under the random 80/20 split (run dated 2026-07-22),
verified via a 14/14 exact match on already-scored equations before
rescoring the originally-null ones (Table~\ref{tab:randomsplit}). This
breaks down as 10 equations passing in the original run, 2 further
equations recovered after patching evaluator defects, 1 near-miss (Planck
blackbody, $\Rsq = 0.999994$, short by $6\times10^{-6}$), 13 genuine
finite/catastrophic extrapolation failures (a separate count from either
headline pass figure --- see the policy note above), and 4 equations left
\textbf{unmeasured} by an architecturally distinct, evaluator-unfixable defect
(Section~\ref{sec:feynman30-defects}, item~4); these 4 are not counted among
the 13 genuine failures above and must not be read as failures.

Under the PCA-directed split (run dated 2026-07-23), \textsc{HypatiaX}
achieves \textbf{13/30 (43.3\,\%)}, verified via a 13/13 exact match on
already-scored equations before rescoring the 17 originally-null ones
(Table~\ref{tab:pcasplit}). This supersedes an 18/30 figure carried in
every prior draft without having been checked against any evaluator
defect. The corrected breakdown is 10 equations already passing plus 3
recovered from the null set (allometric scaling law, Snell's law, ideal
gas law), 2 near-misses (Henderson--Hasselbalch, $\Rsq = 0.9998196$; Nernst
equation, $\Rsq = 0.9999372$), 15 genuine finite/catastrophic failures, and
0 unresolved cases --- the placeholder-token defect that blocks 4
random-split equations was checked directly against every PCA-split
formula string and confirmed absent. The previously reported 142/180
figure is likewise withdrawn: it is the exact arithmetic doubling of a
single 90-record run, produced by pooling two same-shaped flat export
files that describe one underlying run, not two distinct 90-instance
pools.

We do not assert a difficulty ranking between the two protocols. The
corrected PCA-split rate (43.3\,\%) exceeds the corrected random-split rate
(40.0\,\%), the opposite of what a ``PCA extrapolation is strictly harder''
narrative would predict. Both figures use the strict $\Rsq \geq 0.999999$
threshold defined in the policy note above and report near-misses
separately; under the near-miss-inclusive convention instead, the
random-split loose tally is $12+1=13/30$ (43.3\,\%) and the PCA-split loose
tally is $13+2=15/30$ (50.0\,\%) --- neither of which should be confused
with the PCA-split \emph{strict} figure of $13/30$ stated above, despite
sharing the same numerator by coincidence. The 4 random-split equations
affected by the placeholder-token defect (Section~\ref{sec:feynman30-defects},
item~4) have no valid $\Rsq$ and are reported as \textbf{unmeasured}, not as
failures. The underlying defect is confirmed architecturally unfixable by
evaluator patching, so unlike the earlier withdrawn figures this is not
treated as a pending correction; the unmeasured status is disclosed as a
permanent limitation of this run, not deferred or silently dropped from the
per-equation table. Table~\ref{tab:feynman30-legacy} is \emph{not} a per-equation
breakdown of the 12/30 or 13/30 headline figures above. It reports a
different, withdrawn 9/30 run (a distinct equation set, hardware
configuration, and scoring pass) and is retained only as a historical
record of that earlier, superseded run.

\begin{longtable}{p{2.3cm}p{6.3cm}p{1.4cm}p{2.5cm}}
\caption{Random 80/20 split, run dated 2026-07-22, all 30 equations, after
  patching evaluator defects and gating on a 14/14 exact-reproduction
  self-check. \textbf{Bold}: pass ($\Rsq \geq 0.999999$).}
\label{tab:randomsplit}\\
\toprule
Domain & Equation & $\Rsq$ & Status \\
\midrule
\endfirsthead
\toprule
Domain & Equation & $\Rsq$ & Status \\
\midrule
\endhead
Biology & Allometric scaling law & \textbf{1.0000} & pass \\
Biology & Logistic growth rate & \textbf{1.0000} & pass \\
Biology & Michaelis--Menten enzyme kinetics & \textbf{1.0000} & pass \\
Chemistry & Arrhenius rate constant & $-2.368$ & fail (finite) \\
Chemistry & Henderson--Hasselbalch buffer pH & \textbf{1.0000} & pass \\
Electrochemistry & Nernst equation & --- & unmeasured$^\dagger$ \\
Electromagnetism & Clausius--Mossotti dielectric field & \textbf{1.0000} & pass \\
Electromagnetism & Dielectric polarisation & $-2.2\times10^{45}$ & catastrophic \\
Electromagnetism & Energy in a capacitor & $-1.240$ & fail (finite) \\
Electromagnetism & Lorentz force & $-\infty$ & catastrophic \\
Electromagnetism & Ohm's law & \textbf{1.0000} & pass \\
Electrostatics & Coulomb force, 1D simplified & $-0.066$ & fail (finite) \\
Electrostatics & Coulomb's law & $-0.253$ & fail (finite) \\
Magnetism & Curie's law & --- & unmeasured$^\dagger$ \\
Mechanics & Kinetic energy (classical) & $-1.590$ & fail (finite) \\
Mechanics & Newton's gravitational force & 0.6634 & fail (finite) \\
Mechanics & Reduced mass, two-body system & \textbf{1.0000} & pass \\
Mechanics & Total mechanical energy & \textbf{0.9999998} & pass (recovered) \\
Optics & Double-slit interference intensity & \textbf{1.0000} & pass \\
Optics & Snell's law & \textbf{1.0000} & pass (recovered) \\
Probability & Gaussian probability density & 0.1340 & fail (finite) \\
Quantum & Bose--Einstein occupation number & --- & unmeasured$^\dagger$ \\
Quantum & Fermi--Dirac occupation number & $-\infty$ & catastrophic \\
Quantum & Photon energy & $-\infty$ & catastrophic \\
Quantum & Rabi frequency & $-3.058$ & fail (finite) \\
Quantum & Zeeman energy & \textbf{1.0000} & pass \\
Thermodynamics & Fourier's law of heat conduction & $-0.2803$ & fail (finite) \\
Thermodynamics & Ideal gas law & \textbf{1.0000} & pass \\
Thermodynamics & Planck blackbody spectral radiance & 0.999994 & near-miss \\
Thermodynamics & Stefan--Boltzmann law & --- & unmeasured$^\dagger$ \\
\midrule
\multicolumn{4}{l}{Pass: 12/30 (40.0\%). Near-miss: 1/30. Genuine fail: 13/30. Unmeasured: 4/30.} \\
\bottomrule
\end{longtable}

{\small $^\dagger$ Placeholder-token defect (Section~\ref{sec:feynman30-defects}, item~4):
stored formula references a variable name that was never a real input, so no
$\Rsq$ could be computed. Reported as \textbf{unmeasured}, not as a failure:
this status must not be read as, or aggregated with, ``fail''.\par}

\begin{longtable}{p{4.8cm}p{1.6cm}p{2.9cm}}
\caption{PCA-directed 40/60 split, run dated 2026-07-23. The 17 rows below
  were originally null and are individually re-derived here, gated on a
  13/13 exact-reproduction self-check against the 13 already-scored
  equations. The remaining 13 already-scored equations (10 pass, 3 fail per
  the aggregate self-verification count) are not itemised here: their
  individual $\Rsq$ values were not captured at per-equation granularity
  during this rescoring run and are not reconstructed by inference.}
\label{tab:pcasplit}\\
\toprule
Equation & $\Rsq$ & Status \\
\midrule
\endfirsthead
\toprule
Equation & $\Rsq$ & Status \\
\midrule
\endhead
Allometric scaling law & \textbf{0.9999995} & pass (recovered) \\
Arrhenius rate constant & $-0.363$ & fail (finite) \\
Henderson--Hasselbalch buffer pH & 0.9998196 & near-miss \\
Nernst equation & 0.9999372 & near-miss \\
Dielectric polarisation & $-1.068$ & fail (finite) \\
Lorentz force & $-1.030$ & fail (finite) \\
Coulomb force, 1D simplified & $-0.863$ & fail (finite) \\
Coulomb's law & $-0.460$ & fail (finite) \\
Snell's law & \textbf{1.0000} & pass (recovered) \\
Gaussian probability density & $-1.192$ & fail (finite) \\
Photon energy & $-13.592$ & catastrophic \\
Bose--Einstein occupation number & $-1.051$ & fail (finite) \\
Rabi frequency & $-13.117$ & catastrophic \\
Planck blackbody spectral radiance & $-27.832$ & catastrophic \\
Fourier's law of heat conduction & $-0.084$ & fail (finite) \\
Stefan--Boltzmann law & $-0.889$ & fail (finite) \\
Ideal gas law & \textbf{0.9999999933} & pass (recovered) \\
\midrule
\multicolumn{3}{l}{Of these 17: 3 pass, 2 near-miss, 12 genuine fail.} \\
\bottomrule
\end{longtable}

\subsubsection*{Confirmed evaluator defects}\label{sec:feynman30-defects}

\begin{enumerate}
\item \texttt{\textasciicircum} parsed as bitwise XOR, not exponentiation.
  Fixed by rewriting to \texttt{**} before evaluation. Random split only.
\item Unstripped bracketed annotation breaks the parser. Fixed by stripping
  a leading \texttt{[\ldots]} segment before evaluation. Random split only.
\item Missing \texttt{safe\_asin}/\texttt{safe\_acos}. Fixed by injecting
  clipped-inverse-trig definitions into the evaluator namespace. Both
  splits.
\item \textbf{Placeholder-token variable names (e.g.\ \texttt{x0},
  \texttt{x1}) baked into stored formulas.} Architecturally distinct from
  items 1--3 and 5: no evaluator patch can recover a reference to a
  variable that was never real, so this is not something a code fix can
  close. Affects 4 random-split equations (Nernst,
  Curie's law, Bose--Einstein, Stefan--Boltzmann); confirmed absent from
  all 30 PCA-split formula strings by direct pattern search. Root cause (an
  upstream formula-serialisation step) is not yet located and is disclosed
  as an open, unresolved limitation rather than a pending fix. These 4
  equations are reported throughout as \textbf{unmeasured} (no valid
  $\Rsq$ was or could be computed) and are excluded from both the pass and
  fail counts; they must not be silently dropped from the per-equation
  table nor read as failures.
\item Missing \texttt{asin\_of\_sin}. Six of the seventeen PCA-split null
  equations failed uniformly with an undefined-name error; confirmed as
  $\texttt{asin\_of\_sin}(x)=\arcsin(\mathrm{clip}(\sin(x),-1,1))$ and
  patched into the evaluator namespace.
\item \texttt{asin\_of\_sin}/\texttt{safe\_asin}/\texttt{acos\_of\_cos}/
  \texttt{safe\_acos} aliased to domain-unsafe bare
  \texttt{arcsin}/\texttt{arccos} by the discovery-time expression
  normaliser in \texttt{hybrid\_system\_v50\_2.py}, silently discarding the
  clip/compose semantics these names exist to provide. Distinct from item
  5: this affects PySR's \emph{internal} RMSE/$\Rsq$ computation during
  formula discovery and candidate selection, not final scoring against the
  production evaluator. Fixed by leaving these names unaliased for
  \texttt{sympy.lambdify} and instead supplying clip/compose-correct
  implementations via its \texttt{modules} argument.
\end{enumerate}

\textbf{Remaining limitations.} One item above is not fully closed: the
placeholder-token root cause (item~4) has not been located, so the 4
affected random-split equations remain permanently unevaluatable. This
does not require a code fix by any deadline --- it is disclosed as a
permanent limitation, not a pending correction --- but the paper must
consistently report these 4 as \textbf{unmeasured} throughout (as done in
Table~\ref{tab:randomsplit} and the surrounding prose above), rather than
omitting them or letting their status be read as ``fail''. Item~5
(\texttt{asin\_of\_sin}) required a correction to how it
was originally reported, and item~6 (discovery-time normalisation) is a
newly-disclosed related defect; both are addressed below.

We obtained and directly inspected the actual production evaluator source
(\texttt{run\_comparative\_suite\_benchmark\_v2.py} and
\texttt{run\_comparative\_suite\_benchmark\_pca.py}, the scripts that
construct the \texttt{safe\_globals} execution namespace and call
\texttt{eval()}/\texttt{exec()} on stored formula strings for scoring ---
as distinct from \texttt{hybrid\_system\_v50\_2.py} and
\texttt{symbolic\_engine.py}, which reference
\texttt{safe\_asin}/\texttt{asin\_of\_sin} only on the PySR
\emph{discovery}/search side, not the scoring side). \textbf{Confirmed:}
all three \texttt{safe\_globals} constructions in
\texttt{run\_comparative\_suite\_benchmark\_v2.py} omitted
\texttt{safe\_asin}, \texttt{safe\_acos}, and \texttt{asin\_of\_sin}
entirely, defining only a bare, clipped \texttt{arcsin}/\texttt{arccos}.
Since PySR can be configured (via \texttt{--use-transcendental-compositions})
to discover formulas that call these exact names --- the script's own CLI
help text states this is ``required to recover equations like Snell's law
at $\Rsq>0.9999$'' --- any such formula raises \texttt{NameError} during
scoring, caught by a broad \texttt{except Exception: pass}, and the
function returns \texttt{None}: a silently null, not a silently wrong,
extrapolation score. This is consistent with, and is the most likely root
cause of, exactly the pattern already observed: \textbf{Snell's law itself
was one of the equations originally null before this paper's rescoring
pass} (Table~\ref{tab:pcasplit}).

\textbf{Correction to a previously-reported fix.} An earlier pass of this
diagnostic claimed the above had been ``fixed by adding \texttt{safe\_asin},
\texttt{safe\_acos}, \texttt{asin\_of\_sin}, \texttt{acos\_of\_cos}, and
\texttt{atan\_of\_tan} to all five \texttt{safe\_globals} constructions in
both production scripts.'' On direct line-by-line inspection this claim
was false as stated: of the three \texttt{safe\_globals} constructions in
\texttt{run\_comparative\_suite\_benchmark\_v2.py}, only one (a module-level
helper used solely by \texttt{compute\_extrap\_r2\_far} for far-region
extrapolation re-scoring) had actually received any patch, and even that
one was incomplete, missing \texttt{acos\_of\_cos} and
\texttt{atan\_of\_tan}. The two \texttt{safe\_globals} constructions that
govern \emph{primary in-range} scoring ---
\texttt{BaseMethod.\_runner\_eval\_formula} (used by the Hybrid,
\texttt{SymbolicEngineMethod}, and \texttt{HybridSystemV50\_2Method}
scoring paths) and \texttt{PureLLMBaselineMethod.\_runner\_eval\_formula}
(the pure-LLM baseline's main scoring path) --- were entirely unpatched, so
the original defect remained live in the paths that actually produce
headline $\Rsq$ figures. This has now been corrected: all three
\texttt{safe\_globals} constructions in
\texttt{run\_comparative\_suite\_benchmark\_v2.py} carry all five names
(\texttt{safe\_asin}, \texttt{safe\_acos}, \texttt{asin\_of\_sin},
\texttt{acos\_of\_cos}, \texttt{atan\_of\_tan}), verified by direct
execution against synthetic calls that previously raised \texttt{NameError}
and now return the correct clipped-composition value. A related,
previously-undisclosed instance of the same defect was also found and
fixed in \texttt{symbolic\_engine.py}'s internal
\texttt{\_predict\_from\_equation} evaluation namespace, which is missing
despite this engine injecting these exact operator names into PySR's
search space --- meaning discovered candidates using them could not even
be scored internally during search, prior to this fix.

\textbf{A second, distinct defect (item~6)} was found in
\texttt{hybrid\_system\_v50\_2.py}: its discovery-time expression
normaliser aliased \texttt{asin\_of\_sin}/\texttt{safe\_asin} to bare
\texttt{arcsin} and \texttt{acos\_of\_cos}/\texttt{safe\_acos} to bare
\texttt{arccos} via literal text substitution before
\texttt{sympy.lambdify}, discarding the clip/compose semantics these names
exist to provide. Concretely, $\texttt{asin\_of\_sin}(x)$ --- designed to
handle $x$ outside $[-1,1]$ by composing through $\sin(\cdot)$ first ---
was silently replaced with plain $\arcsin(x)$, reproducing exactly the
domain failure it was designed to avoid and yielding \texttt{NaN} for
out-of-domain inputs. This is architecturally distinct from item~5: it
affects PySR's \emph{internal} RMSE/$\Rsq$ computation used for candidate
fitness and selection \emph{during} formula discovery, not the final
scoring reported in this paper's tables. It has been fixed by leaving
these names unaliased through normalisation and instead supplying
clip/compose-correct implementations to \texttt{sympy.lambdify} via its
\texttt{modules} argument, verified to return finite values (rather than
\texttt{NaN}) for out-of-domain test inputs where the unfixed alias-based
path did not.

\textbf{What remains unverified.} This corrected pass directly inspected
\texttt{run\_comparative\_suite\_benchmark\_v2.py},
\texttt{run\_comparative\_suite\_benchmark\_pca.py},
\texttt{symbolic\_engine.py}, and \texttt{hybrid\_system\_v50\_2.py}.
\texttt{run\_comparative\_suite\_benchmark\_pca.py} was found to already
carry all five names (\texttt{safe\_asin}, \texttt{safe\_acos},
\texttt{asin\_of\_sin}, \texttt{acos\_of\_cos}, \texttt{atan\_of\_tan}) in
all three of its \texttt{safe\_globals} constructions, with correct
clipped/composed semantics; this was confirmed not merely by reading the
source but by extracting each \texttt{safe\_globals} dictionary literal
and executing it directly against out-of-domain test inputs (e.g.\
$\texttt{asin\_of\_sin}(10.0)$), all of which returned finite,
correctly-clipped values rather than raising or returning \texttt{NaN}.
Given that the equivalent claim for
\texttt{run\_comparative\_suite\_benchmark\_v2.py} was demonstrably false
on first inspection, this PCA-split result should not be read as a
foregone conclusion re-confirmed by analogy --- it was checked
independently and holds up. Item~5 and item~6 are therefore now closed for
all four files inspected. Separately, none of the fixes described above
have yet been exercised against a full production rerun of either
benchmark script, and the item~6 fix means any past discovery run that
used these five operator names may have selected or reported formulas
based on an internal fitness score computed against the wrong
(bare-\texttt{arcsin}) function --- a possible influence on which formulas
were surfaced at all, not only how they were subsequently scored. Until a
full rerun is obtained, all of the above should be treated as
applied-in-code and verified-in-isolation only, on the same footing as
other applied-but-not-yet-reconfirmed-in-production fixes elsewhere in
this paper. None of this changes the 12/30 or 13/30 figures reported
above, which are already gated on exact reproduction of every
previously-scored equation and already used a manually-patched evaluator
copy with the item~5 fix in place; but a full production rerun could
reveal additional previously-null equations (beyond Snell's law) now
scoring successfully, in either split, which this diagnostic pass ---
confined to the already-identified null set --- would not have surfaced.

\textbf{Historical note: the withdrawn 9/30 run.}\quad
Table~\ref{tab:feynman30-legacy} reports the earlier, withdrawn Kaggle
4-vCPU run on its own (different) equation set, together with that run's
own Mann-Whitney statistics and hardware-sensitivity comparison; it is
retained only as a historical record and is not the source of, or a
per-equation breakdown for, the 12/30 or 13/30 figures above.

\begin{table}[h]
\centering
\caption{\textbf{[WITHDRAWN RUN --- NOT THE SOURCE OF THE PAPER'S HEADLINE FIGURES]}
  Feynman extrapolation benchmark, Kaggle 4-vCPU multiprocessing run,
  9/30 solve rate, withdrawn (\S\ref{sec:feynman30}).
  \textbf{Bold}: extrap $\Rsq > 0.99$; \emph{italic}: $\Rsq < 0$.
  This table uses a different equation ordering and hardware configuration
  than the corrected runs and is retained only as a historical record. The
  paper's reported results are \textbf{12/30 (40.0\,\%)} successes under the
  random 80/20 split (2026-07-22) and \textbf{13/30 (43.3\,\%)} under the
  PCA-directed 40\,\%/60\,\% split (2026-07-23); no per-equation table for
  either of those two runs currently exists. Do not cite the 9/30 total
  below, or any row in this table, as support for the 12/30 or 13/30 figures.}
\label{tab:feynman30-legacy}
\small\setlength{\tabcolsep}{4.5pt}
\renewcommand{\arraystretch}{1.08}
\begin{tabular}{llccc}
\toprule
\textbf{Equation} & \textbf{Domain} &
\textbf{Hyp Train $\Rsq$} & \textbf{Hyp Extrap $\Rsq$} & \textbf{NN Extrap $\Rsq$} \\
\midrule
Gaussian                  & Mechanics      & 0.926 & $-10.36$           & $-24.20$ \\
Coulomb Force             & Mechanics      & 0.869 & $-7.43$            & $\ll{-}100$ \\
Relativistic momentum     & Mechanics      & 0.997 & $-0.25$            & $-4.76$ \\
Doppler shift             & Mechanics      & 0.997 & 0.688              & $-0.26$ \\
Harmonic oscillator       & Mechanics      & 0.997 & \textbf{1.000}     & $-2.04$ \\
Electric potential energy & Thermodynamics & 0.997 & \textbf{0.998}     & 0.962 \\
Energy of photon          & Thermodynamics & $-2.71$ & $\ll{-}100$      & 0.677 \\
Magnetization             & Thermodynamics & 0.924 & $-0.47$            & $-1.07$ \\
Relativistic Doppler      & Optics         & 0.998 & \textbf{0.994}     & 0.987 \\
Heat conduction           & Optics         & 0.923 & 0.136              & $\ll{-}100$ \\
Snell's law               & Optics         & 0.993 & $-0.31$            & $-0.13$ \\
Polarization              & Electromagnetism & 0.982 & 0.941            & 0.923 \\
Torque                    & Electromagnetism & 0.998 & \textbf{1.000}   & $-2.01$ \\
Interference intensity    & Electromagnetism & 0.985 & \textbf{1.000}   & $-6.07$ \\
Polarizability            & Electromagnetism & $-0.95$ & $-11.75$       & 0.931 \\
Planck radiation          & Electromagnetism & $-0.86$ & $-5.90$        & $-1.39$ \\
Photon energy             & Quantum        & $-2.61$ & $\ll{-}100$      & 0.906 \\
Magnetic moment           & Quantum        & $-0.76$ & $-9.59$          & $-2.56$ \\
Bose-Einstein             & Quantum        & 0.997 & \textbf{0.997}     & 0.778 \\
Gravity potential         & Gravitation    & 0.978 & $-2.38$            & $\ll{-}100$ \\
Orbital period            & Gravitation    & 0.998 & \textbf{1.000}     & 0.862 \\
Dielectric constant       & Fluid          & 0.579 & 0.000              & 0.000 \\
Diffraction               & Fluid          & 0.995 & \textbf{0.997}     & 0.825 \\
Wave superposition        & Waves          & 0.692 & $-1.14$            & $\ll{-}100$ \\
de Broglie wavelength     & Waves          & $-0.11$ & $-9.46$          & $\ll{-}100$ \\
Time dilation             & Relativity     & 0.997 & 0.639              & $-1.78$ \\
Lorentz factor            & Relativity     & 0.997 & 0.711              & $-0.54$ \\
Coulomb potential         & Atomic         & 0.063 & $\ll{-}100$        & $-4.66$ \\
Diffusion coefficient     & Atomic         & $-0.56$ & $\ll{-}100$      & 0.034 \\
Larmor frequency          & Nuclear        & 0.998 & \textbf{1.000}     & $-1.40$ \\
\midrule\midrule
\textbf{Successes ($\Rsq>0.99$), this run (withdrawn, \S\ref{sec:feynman30})} & & & \textbf{9/30 (30.0\%)} & 0/30 (0.0\%) \\
\textbf{Median extrap $\Rsq$}    & & & $-0.123$               & $-1.225$ \\
MW all $n=30$       & \multicolumn{4}{l}{$U=534.5$, $p=0.107$ — not significant (expected)} \\
\textbf{MW successes $n=9$} & \multicolumn{4}{l}{$U=81$, $p=\mathbf{0.00020}$ — strongly significant} \\
\bottomrule
\end{tabular}
\end{table}

\emph{The remainder of this note describes the withdrawn 9/30 run only and
is not a description of the canonical 12/30 (2026-07-22) or 13/30
(2026-07-23) results above.} The full-benchmark Mann-Whitney test
($U=534.5$, $p=0.107$) on that withdrawn run is not significant, as
expected: 21 failures introduce high variance that masks the signal.
Among the 9 success-equations of that run where \textsc{HypatiaX} achieved
$\Rsq > 0.99$, the conditional Mann-Whitney test ($U=81$, $p=0.00020$) was
strongly significant. Five of those 9 successes corresponded to exact
symbolic recovery ($\Rsq = 1.000$). These statistics have not been
recomputed against the canonical July~22/23 runs.

\textbf{Hardware sensitivity (withdrawn run).}  Across that same withdrawn
benchmark, three independent hardware configurations yielded 14, 8, and 9
successes (Celeron multi-threaded, Colab T4 serial, Kaggle 4-vCPU
multiprocessing respectively), a variation attributed to non-determinism in
the PySR evolutionary search under different parallelism regimes. All three
counts belong to the withdrawn run and have not been recomputed under the
corrected protocol; the $p=0.00020$ and $p=0.00043$ Mann-Whitney results
quoted for the Kaggle and Colab configurations likewise describe only that
withdrawn run.

\textbf{Hardware sensitivity (current pipeline, CI-verified).} The ad hoc
local Celeron/Colab/Kaggle comparison above has been superseded by a
GitHub Actions CI pipeline (repository \texttt{sednabcn/LLM-HypatiaX-REPRO}),
which runs the canonical benchmark on a single, standard
\texttt{ubuntu-24.04} GitHub-hosted runner with \texttt{PySR\_SEED=42} fixed
throughout, i.e.\ one deterministic machine rather than three divergent
ones. This removes hardware sensitivity as a source of variation \emph{in
principle}, but does not by itself confirm the fixes described in
Section~\ref{sec:feynman30-defects} (items~5--6) are reflected in the
pipeline CI is currently exercising. On direct inspection of the CI logs:
the \texttt{exp2\_feynman} discovery job (run~30258093290, commit
\texttt{d5adf4e}, 2026-07-27) completed cleanly, and the subsequent
\texttt{exp2\_feynman\_extrap} scoring job (run~30329546386, commit
\texttt{c694d5f}, 2026-07-28) shows \emph{partial, inconsistent} recovery
on exactly the equations items~5--6 concern: Snell's law and the Nernst
equation return a valid (non-null) extrapolation $\Rsq$ for 2 of 6 methods
(\texttt{SymbolicEngineWithLLM}, \texttt{HybridDiscoverySystem v50\_2}) but
remain \texttt{null} for the other 4 (\texttt{PureLLM Baseline},
\texttt{ImprovedNN}, \texttt{EnhancedHybridSystemDeFi},
\texttt{HybridSystemLLMNN}), while the placeholder-token-affected
Bose--Einstein and Stefan--Boltzmann equations remain \texttt{null} across
all methods, consistent with item~4 being architecturally unfixable. This
indicates the evaluator-namespace fixes for items~5--6 have not yet
propagated to whatever commit this CI pipeline is currently running against.
\textbf{Item~9 is therefore closed with respect to its original scope} ---
the stale local hardware-variance numbers are no longer the operative
reproducibility story, and are superseded by a citable, deterministic,
single-runner CI pipeline --- \textbf{but this CI evidence does not yet
support quoting a corrected 12/30- or 13/30-style headline number}; it
instead corroborates, on live production infrastructure, that the item~5--6
defects were real and are not yet fully resolved in whatever code CI is
running. A rerun after confirming the evaluator patches are on the commit
CI builds from is needed before any new headline figure can be cited.

% ── RF-04: Nguyen-12 Benchmark ───────────────────────────────────────────────
\subsection{Nguyen-12 Benchmark (Standard SR Suite)}\label{subsec:nguyen12-benchmark-standard-sr-suite}
\label{sec:nguyen12}

The Nguyen-12 suite~\citep{uy2011semantically} is the standard symbolic-regression
benchmark comprising 12 polynomial and transcendental equations (Nguyen-1 through
Nguyen-12) with up to two input variables.  We evaluate HypatiaX and PySR-only on
the same PCA-directed aggressive extrapolation protocol used for the DeFi benchmark
(40\,\%/60\,\% split), and compare against the neural MLP baseline.

\begin{table}[h]
\centering
\caption{Nguyen-12 benchmark: train and extrapolation $\Rsq$ by equation.
$P$ = PySR-only; $H$ = HypatiaX; $N$ = Neural MLP.
Near-miss criterion: $\Rsq \ge 0.9999$ (4-decimal convention of
\citealt{uy2011semantically}).
Near-miss detail: $\Rsq\in[0.9998,0.9999)$ with perfect extrapolation
($\Rsq=1.0000$ on $2\times$ holdout). Strict recovery ($\ge 0.9999$) yields
4/12 (33.3\%); the 91.7\,\% figure uses the 4-decimal rounding convention.}
\label{tab:nguyen12}
\begin{tabular}{llrrrrrr}
\toprule
Eq. & Formula &
  \multicolumn{2}{c}{PySR-only ($P$)} &
  \multicolumn{2}{c}{HypatiaX ($H$)} &
  \multicolumn{2}{c}{Neural MLP ($N$)} \\
\cmidrule(lr){3-4}\cmidrule(lr){5-6}\cmidrule(lr){7-8}
 &  & Train & Extrap & Train & Extrap & Train & Extrap \\
\midrule
N-1  & $x^3 + x^2 + x$
  & 0.9999 & \textbf{1.0000}
  & 0.9999 & \textbf{0.9999}
  & 0.9993 & $-0.784$ \\
N-2  & $x^4 + x^3 + x^2 + x$
  & 0.9999 & \textbf{1.0000}
  & 0.9999 & \textbf{1.0000}
  & 0.9986 & $-0.902$ \\
N-3  & $x^5 + x^4 + x^3 + x^2 + x$
  & 0.9999 & \emph{$-426.2$}
  & 0.9999 & \textbf{0.9976}
  & 0.9986 & $-0.913$ \\
N-4  & $x^6 + x^5 + x^4 + x^3 + x^2 + x$
  & 0.9999 & \emph{$\ll{-}100$}
  & 0.9999 & \emph{$\ll{-}100$}
  & 0.9979 & $-0.828$ \\
N-5  & $\sin(x^2)\cos(x)-1$
  & 0.9999 & \textbf{1.0000}
  & 0.9999 & \textbf{1.0000}
  & 0.9979 & $-5.586$ \\
N-6  & $\sin(x)+\sin(x+x^2)$
  & 0.9999 & \textbf{1.0000}
  & 0.9999 & \textbf{1.0000}
  & 0.9987 & $-12.654$ \\
N-7  & $\ln(x+1)+\ln(x^2+1)$
  & 0.9999 & 0.9762
  & 0.9999 & 0.7316
  & 0.9868 & $0.856$ \\
N-8  & $\sqrt{x}$
  & 0.9999 & \textbf{1.0000}
  & 0.9999 & \textbf{1.0000}$^\dag$
  & 0.9988 & $0.954$ \\
N-9  & $\sin(x)+\sin(y^2)$
  & 0.9999 & \textbf{1.0000}
  & 0.9999 & \textbf{1.0000}
  & 0.9986 & $-6.708$ \\
N-10 & $2\sin(x)\cos(y)$
  & 0.9999 & \textbf{1.0000}
  & 0.9999 & \textbf{0.9997}
  & 0.9995 & $-2.379$ \\
N-11 & $x^y$
  & 0.9999 & \textbf{1.0000}
  & 0.9999 & \textbf{0.9999}
  & 0.9984 & $-0.423$ \\
N-12 & $x^4 - x^3 + \tfrac{1}{2}y^2 - y$
  & 0.9987 & $-1.056$
  & 0.9994 & $-1.054$
  & 0.9985 & $-1.198$ \\
\midrule
\multicolumn{2}{l}{Success ($\Rsq\ge0.9999$)}
  & \multicolumn{2}{c}{10/12 (83.3\%)}
  & \multicolumn{2}{c}{\textbf{11/12 (91.7\%)}}
  & \multicolumn{2}{c}{0/12 (0\%)} \\
MW $H>P$   & \multicolumn{7}{l}{$U=51.0$, $p=0.893$ — not significant} \\
MW $P>N$   & \multicolumn{7}{l}{$U=113.0$, $p=0.0097$ — significant} \\
\bottomrule
\end{tabular}
\begin{tablenotes}\small
  \item \textbf{Bold}: extrap $\Rsq \ge 0.9999$. \emph{Italic}: $\Rsq < 0$.
    $P$ = PySR-only; $H$ = HypatiaX; $N$ = Neural MLP.
  \item[$\dag$] N-8 ($\sqrt{x}$): HypatiaX routed to \texttt{hybrid\_llm\_only} mode.
    The LLM directly proposed $x^{0.5}$, which passed the trust gate with
    train $\Rsq = 0.9999$ and achieved exact extrapolation ($\Rsq = 1.000$)
    without invoking PySR. Wall-clock time: 7.4\,s vs.\ 301\,s for PySR-only ---
    the only equation where the LLM warm-start fully bypassed evolutionary search.
\end{tablenotes}
\end{table}

\paragraph{Statistical test.}
Mann-Whitney U-test comparing SR (PySR-only + HypatiaX combined) vs.\ Neural MLP
extrapolation $\Rsq$ across all 12 equations: $U=113$, $p=0.0097$ (two-sided),
confirming that symbolic regression significantly outperforms the neural baseline
on this standard suite under aggressive extrapolation.

\subsection{Stability Under Stochastic Inference}\label{subsec:stability-under-stochastic-inference}
\label{sec:instability_results}

To assess LLM instability systematically, we run $K=30$ independent stochastic
evaluations per task.  The Instability Index $\II_i = \sigma_i$ (standard deviation
of $\Rsq$ across runs) is reported for each task.

Table~\ref{tab:instability} summarises the regime distribution across the 70 tasks
for which multi-run data are available.

\begin{table}[h]
\centering
\caption{LLM instability regime distribution (70 tasks, $K=30$ runs each).
$\II_i = \sigma_i = \mathrm{std}(R^2_i)$ across independent runs.}
\label{tab:instability}
\begin{tabular}{lrrrr}
\toprule
Regime & Definition & $n$ & Fraction \\
\midrule
A: Symbolic Stability   & $\sigma\approx0$, $\mu\approx1$            & 61 & 87.1\,\% \\
B: Deterministic Biased & $\sigma\approx0$, $\mu<1$                  &  2 & 2.9\,\%  \\
B*: Borderline Stochastic & $0 < \sigma < 0.05$                      &  4 & 5.7\,\%  \\
C: Stochastic Collapse  & $\sigma \ge 0.10$ or $\mu < 0$             &  3 & 4.3\,\%  \\
\bottomrule
\end{tabular}
\end{table}

The distribution is strongly bimodal: 87\,\% of tasks exhibit perfect symbolic
stability ($\II=0$), while a small but critical tail of 3 tasks undergoes full
stochastic collapse.

\textbf{Qualification on C-Collapse regime.}  The C-Collapse regime
(Regime~C) comprises only 2--3 cases in the current benchmark; regime-level
claims should be treated as \emph{preliminary evidence} pending a larger
C-Collapse sample.  A reviewer should not read the Spearman $\rho=-0.70$ result
as driven by the C-Collapse cases alone: the correlation is computed across all
70 tasks, with A-Symbolic cases at ($\II=0$, $\Rsq=1$) and C-Collapse cases at
(high $\II$, low $\Rsq$) forming the two poles of the distribution.  The
\textbf{aggregate Spearman $\rho=-0.70$ ($p<0.001$) is the primary statistical
result}; the regime typology is an interpretive overlay.

\textbf{Note on OOD proxy.}  The ``out-of-distribution $\Rsq$'' used to compute
the Spearman correlation is a \emph{test-split proxy} (eval\_mode=test\_r2\_proxy),
not true symbolically-verified OOD performance.  See Section~\ref{sec:ood_proxy}
for full disclosure and the planned Mode~A upgrade.  The proxy likely
\emph{underestimates} the instability--extrapolation gap; we expect $\rho$ to
strengthen under true OOD evaluation.

One C-Collapse case (Portfolio Expected Shortfall) achieves proxy OOD $\Rsq=1.000$
despite $\II=0.311$.  This anomaly requires true OOD evaluation (Mode~A:
\textsc{SymPy} + \texttt{generate\_data} sampling beyond the training range) to
resolve.  If Mode~A reclassifies this case to A-Symbolic, the C-Collapse regime
reduces to $n=2$ and the regime-level claim becomes further preliminary.

The three Regime C tasks are:
\begin{itemize}
  \item \textbf{Portfolio Expected Shortfall for correlated assets}:
    $\mu=0.709$, $\II=0.311$, $p(>0.9)=53\,\%$.  High instability combined with
    moderate mean accuracy suggests that the model alternates between a correct
    multi-asset correlation formula and an incorrect single-asset approximation.
  \item \textbf{Theta of option}: $\mu=0.569$, $\II=0.148$, $p(>0.9)=3\,\%$.
    Only 13 of 30 runs completed without NaN output; the others produced NaN due
    to division by zero in the partial derivative evaluation.
  \item \textbf{Black--Scholes Call Price}: $\mu=0.229$, $\II=0.267$, $p(>0.9)=0\,\%$.
    The most severe failure: zero successful runs in 28 attempts.  The model
    appears to confuse alternate parameterisations of $d_1$ and $d_2$ across runs,
    producing structurally plausible but numerically incorrect formulas.
\end{itemize}

% =============================================================================
\section{Discussion}
\label{sec:discussion}

\subsection{Three Core Findings}\label{subsec:three-core-findings}
\label{sec:disc_findings}

Our results establish three core findings.

\paragraph{Finding 1: Pure LLM approaches are unreliable.}
The bimodal performance of the LLM baseline (median $\Rsq=1.00$, mean $\Rsq=-0.76$)
confirms that a single-run $\Rsq$ measurement is a misleading indicator of LLM
reliability.  Six catastrophic failures in 74 tasks (8.1\,\%) represent an
unacceptable failure rate for production deployment in a financial system.  The
Instability Index makes this failure mode explicit: a task with $\II>0.10$ has
zero probability of near-perfect recovery.

\paragraph{Finding 2: Neural networks fail systematically under extrapolation.}
The neural MLP achieves median $\Rsq=-0.47$ on our benchmark, confirming that
standard deep learning is not a solution to the extrapolation problem in DeFi
mathematics.  Even with data augmentation and residual architecture, the MLP does
not recover the functional form of the underlying formula.  This is not a
hyperparameter tuning failure; it is a fundamental architectural limitation under
distribution shift.

\paragraph{Finding 3: Reliability is the defining advantage of hybrid systems.}
HypatiaX's most important result is not the 89.2\,\% near-perfect success rate
\textbf{[flagged: corrected 60.8\,\%, \S\ref{sec:hybrid-attribution-bug}]}
(impressive as that is) but the \textbf{elimination of catastrophic failures}.
A system that achieves 90\,\% perfect recovery but fails catastrophically 10\,\%
of the time is \emph{less} useful in production than a system that achieves 80\,\%
perfect recovery but fails gracefully on the remainder.  HypatiaX achieves both:
high peak performance and zero catastrophic failures. This finding is precisely
why the bug quantified in \S\ref{sec:hybrid-attribution-bug} does not undermine
the paper's central thesis: catastrophic-failure elimination, not the exact
near-perfect percentage, is the load-bearing claim, and that claim is
unaffected by the correction.

\subsection{Arrhenius: A Verified Rescue Case}\label{subsec:arrhenius-and-scaleincompatibility-failu}
\label{sec:arrhenius_failure}

This subsection was previously withdrawn pending re-derivation from real
data; it is restored here with the verified values from
Table~\ref{tab:llm_ablation}. The Arrhenius rate law $k = A\,e^{-E_a / RT}$
is a genuine rescue case, not a shared failure. PySR-only collapses at
$5\times$ training range (far-$\Rsq = -12.5552$), while HypatiaX recovers
the correct functional form and extrapolates successfully (far-$\Rsq =
0.9028$), despite HypatiaX's training $\Rsq$ (0.9979) being only marginally
higher than PySR-only's (0.9974) --- the far-extrapolation gap is far larger
than the training-fit gap would suggest. This is consistent with the
routing cascade's scale-compatibility check (Figure~\ref{fig:routing_cascade},
annotated ``prevents Arrhenius''): the LLM-warm-started search appears to
avoid a local minimum that the cold PySR-only search falls into on this
equation. We do not have a confirmed mechanistic account of \emph{why} the
warm start avoids this particular local minimum, and do not claim more than
the observed outcome.

We do not extend this explanation to the Portfolio Variance seed-sweep
degradation (Section~\ref{sec:portfolio_seed_sweep}): that case involves a
different failure pattern (HypatiaX underperforming PySR-only at seed~99
specifically) and has not been independently traced to the same or a
different mechanism. Treat the two as separate, unexplained phenomena
rather than instances of one shared cause.

\subsection{The Stability--Accuracy Distinction}\label{subsec:the-stabilityaccuracy-distinction}
\label{sec:disc_stability}

Standard machine learning evaluation collapses the reliability question into a
single point estimate: the test accuracy on a fixed split.  This is insufficient
for systems that will be deployed across a distribution of tasks with varying
complexity.  We argue that three dimensions of evaluation are necessary:
\begin{enumerate}[leftmargin=2em,label=(\roman*)]
  \item \textbf{Accuracy}: what is the expected performance on a typical task?
  \item \textbf{Reliability}: what is the worst-case performance?
    (Catastrophic failure rate, minimum $\Rsq$ across tasks.)
  \item \textbf{Stability}: across repeated runs on the same task, how consistent
    is the output?  (Instability Index $\II_i$.)
\end{enumerate}

HypatiaX addresses all three dimensions, whereas the pure LLM and neural baselines
each address at most one.

\subsection{Why Transcendental Formulas Fail}\label{subsec:why-transcendental-formulas-fail}
\label{sec:disc_transcendental}

The Gaussian CDF $\Phi$ is the proximate cause of the three Regime C collapses.
Unlike polynomial or exponential expressions, $\Phi$ has no finite closed-form
representation in elementary operations.  Multiple valid approximations exist:
$\Phi(x) \approx (1 + e^{-1.7x})^{-1}$, error-function series, and Taylor
expansions.  Under temperature sampling, the LLM produces different approximations
on different runs---each locally plausible, but producing different numerical
predictions.

This observation suggests a practical mitigation: for formulas containing $\Phi$,
provide the LLM with a canonical implementation in the prompt rather than asking it
to generate one.  We leave this prompt-engineering intervention for future work.

\subsection{OOD Metric as a Proxy}\label{subsec:ood-metric-as-a-proxy}
\label{sec:ood_proxy}

Throughout this paper, ``out-of-distribution $\Rsq$'' refers to performance on a
held-out \emph{test split} constructed by PCA-directed ordering, not to a
symbolically verified out-of-distribution regime.  This distinction matters: a
formula that achieves far-$\Rsq \approx 0$ on the test split may still be
structurally wrong.  The clearest illustration is \textsc{HypatiaX} seed=42 on
Portfolio Variance: far-$\Rsq = -0.023$ (a near-miss on the metric) but the
recovered expression \texttt{(s1+s2)*(rho*0.243+0.738)} is a linear approximation
that is structurally incorrect.  A test-split proxy cannot detect this; symbolic
verification via \texttt{sympy} OOD sampling would catch it immediately.

A planned upgrade (Mode~A symbolic verification) will: (1) add \texttt{formula} and
\texttt{var\_ranges} fields to the results JSON; (2) enable \texttt{sympy} OOD
sampling over $[5\times, 10\times, 100\times]$ the training range; and (3) replace
the \texttt{test\_r2\_proxy} column with a verified \texttt{ood\_sympy} metric.
This is a concrete, numbered future-work item.  The current proxy likely
\emph{underestimates} the instability--extrapolation gap: a structurally incorrect
formula that happens to score well on the held-out test split (same distribution)
will be reclassified as a failure under true OOD evaluation.  We therefore
expect the Spearman $\rho=-0.70$ result to \emph{strengthen} under Mode~A.

\subsection{Design Principles for Analytical Discovery}\label{subsec:design-principles-for-analytical-discove}
\label{sec:disc_principles}

Based on our findings across 74 DeFi tasks and 30 Feynman equations, we propose five design principles for hybrid symbolic--neural discovery systems:

\begin{enumerate}
  \item \textbf{Cascade validation rather than single-criterion acceptance.} Our five-stage routing cascade detected failures missed by any individual gate. Dimensional consistency, domain constraints, numerical stability, OOD probing, and trust scoring each catch distinct failure modes.
  \item \textbf{Multi-scale extrapolation testing.} Evaluate candidates at three extrapolation distances ($1.2\times$, canonical, $5\times$) to distinguish local approximations from globally correct forms. Extrapolation errors increase exponentially with distance from the training range.
  \item \textbf{Quantify instability, not just accuracy.} A single-run $\Rsq$ is a misleading indicator for transcendental formulas. The Instability Index ($\II_i = \sigma_i$ across $K=30$ runs) identifies high-risk tasks before deployment.
  \item \textbf{Domain-aware routing.} LLM guidance is reliable for algebraic and well-documented equations; transcendental token detection (Stage~2) correctly flags cases where symbolic search should dominate.
  \item \textbf{Honest runtime reporting.} Aggregate speedup statistics conflate fast LLM-routed cases with slow neural-fallback cases. Report speedup separately for each routing path.
\end{enumerate}

\subsection{Comparison with Related Work}\label{subsec:comparison-with-related-work}
\label{sec:disc_related}

HypatiaX builds on three lines of prior work. AI Feynman~\citep{udrescu2020ai} uses neural networks to identify symmetries and simplifications before symbolic regression, achieving strong performance on physics equations. Physics-informed neural networks~\citep{raissi2019physics,karniadakis2021physics} offer an alternative by embedding physical constraints directly into the network architecture; our approach differs by targeting financial and biochemical domains where such constraints are unavailable, and by recovering interpretable closed-form expressions rather than learned surrogates. DSR~\citep{petersen2021deep,landajuela2022unified} uses reinforcement learning and risk-seeking policy gradients to guide symbolic search; our LLM warm-start provides a complementary initialisation strategy with lower computational cost. MCTS-based SR methods~\citep{sun2022symbolic} explore expression trees more systematically than evolutionary search; HypatiaX's five-stage routing provides a practical approximation with lower overhead. On the Feynman benchmark, HypatiaX achieves 12/30 successes (40.0\%) under a random 80/20 extrapolation protocol and 13/30 (43.3\%) under our aggressive PCA-directed extrapolation protocol (both at the single strict $\Rsq \geq 0.999999$ threshold defined in \S\ref{sec:feynman30}), comparable to AI Feynman 2.0 under similar conditions. The DeFi benchmark is, to our knowledge, the first symbolic regression evaluation on post-2020 financial mathematics.

\subsection{Ethical Considerations}\label{subsec:ethical-considerations}
\label{sec:disc_ethics}

Automated formula discovery carries risks in high-stakes domains. A formula with high training $\Rsq$ may fail catastrophically under distribution shift---precisely the scenario our benchmark targets. We make three recommendations: (1) never deploy a discovered formula in production without out-of-distribution validation; (2) report instability alongside accuracy; (3) treat HypatiaX output as a hypothesis to be verified, not a certified result. The DeFi domain carries additional risks: discovered liquidation or VaR formulas used in production could amplify systemic risk if they fail under market stress. Our validation framework catches many errors, but not all---particularly the silent semantic failures (correct structure, wrong constants) that produce $\Rsq \approx 1$ on training data.

\subsection{Broader Applicability}\label{subsec:broader-applicability}
\label{sec:disc_broader}

The hybrid symbolic--neural architecture generalizes beyond DeFi to any domain where: (1) the underlying relationship admits a closed-form expression; (2) LLM training data includes related equations; and (3) extrapolation beyond the training range is required. Candidate domains include pharmacokinetics (drug absorption/elimination kinetics), materials science (stress-strain relationships, diffusion coefficients), climate science (radiative forcing equations), and epidemiology (compartmental model parameters). The instability analysis framework applies wherever stochastic inference is used---including neural architecture search and reinforcement learning policy optimization.

\subsection{Limitations and Future Work}\label{subsec:limitations-and-future-work}
\label{sec:limitations}

\paragraph{Single LLM.}
All results use a single LLM (not disclosed for double-blind review).
The degree to which these findings generalise to other LLMs---particularly
models with specialised mathematical pre-training---is an open question.

\paragraph{Benchmark scope.}
The benchmark covers DeFi-specific formulas.  Generalisation to broader scientific
or engineering domains requires additional validation.

\paragraph{Runtime on fallback cases.}
The 120-second NN training cap is a practical constraint that may not reflect the
optimal training time for each task.  A lighter neural architecture with faster
training could improve the aggregate mean time substantially.

\paragraph{Prompt design.}
The trust gate and routing decisions are sensitive to prompt quality.  A systematic
study of prompt design as a hyperparameter is needed before deploying HypatiaX
in new domains.

% =============================================================================
\section{Conclusion}
\label{sec:conclusion}

We have presented HypatiaX, a hybrid symbolic--neural framework for extrapolation
in DeFi mathematical discovery.  On a benchmark of 74 DeFi tasks with an aggressive
PCA-directed extrapolation split, HypatiaX achieves:

\begin{itemize}
  \item \textbf{89.2\,\%} near-perfect success rate ($\Rsq > 0.99$), up from
    62.2\,\% for the pure LLM and 5.4\,\% for the neural baseline
    \textbf{[flagged: corrected figure is 60.8\,\%, which is \emph{below}
    the pure-LLM baseline of 62.2\,\%, not above it --- see
    \S\ref{sec:hybrid-attribution-bug}]}.
  \item \textbf{Zero catastrophic failures}, eliminating the 6 failures ($\Rsq < -10$)
    of the pure LLM baseline (unaffected by the correction above).
  \item \textbf{Performance gains that scale with difficulty}: $+12.5$\,pp on easy tasks,
    $+31.1$\,pp on medium, and $+38.1$\,pp on hard transcendental tasks
    \textbf{[flagged: corrected gains are $0.0$/$0.0$/$-4.8$\,pp respectively---
    HypatiaX ties the pure-LLM baseline on Easy and Medium and underperforms it
    on Hard once decision-attribution is corrected for; the difficulty-scaling
    trend does not survive correction, see \S\ref{sec:hybrid-attribution-bug}]}.
  \item \textbf{$1.73\times$ speedup} over the neural baseline on LLM-routed cases
    ($n=68$), with the caveat that aggregate mean time is higher when fallback NN
    retraining is included.
\end{itemize}

These results establish that hybrid symbolic--neural systems are not merely
competitive with their individual components: they achieve qualitatively different
reliability properties that neither component can achieve alone.  The elimination
of catastrophic failures, in particular, is a prerequisite for deployment in
financial systems where a single incorrect formula can cause significant harm.

Future work should explore: extending the benchmark to broader scientific domains;
developing lighter fallback architectures to reduce mean wall-clock time;
integrating iterative LLM self-correction as an intermediate routing option between
pure symbolic and neural fallback; and systematic evaluation across multiple LLM
families.

% =============================================================================
\section*{Reproducibility Statement}

All results are generated using HypatiaX v3.0 with the following artefacts:
\begin{itemize}
  \item \texttt{hypatiax\_defi\_benchmark\_v3\_results\_m27.json}: full per-task
    results including $\Rsq$, timing, and routing decisions.
  \item \texttt{portfolio\_variance\_seed\_sweep.json}: raw per-seed results for
    the 5-seed Portfolio Variance robustness study (Section~\ref{sec:portfolio_seed_sweep}).
  \item Fixed random seed 42 for all neural network training.
  \item Deterministic prompting framework for LLM formula generation.
  \item Unified formula executor (single code path for all methods).
  \item Checkpointing and resume functionality for long-running experiments.
\end{itemize}
The pipeline includes explicit corrections for data leakage, split misconfiguration,
formula evaluation inconsistency, and trust-gating threshold, as described in
Section~\ref{sec:corrections}.

\textbf{Hardware sensitivity (Feynman benchmark, withdrawn run).}  The Feynman $n=30$ experiment was
run on three hardware configurations: Celeron multi-threaded (14/30 successes),
Colab T4 GPU serial (8/30), and Kaggle 4-vCPU multiprocessing (9/30, primary).
The variation arises from non-determinism in PySR's evolutionary search under
different parallelism regimes. \textbf{Author note:} all three figures belong
to the withdrawn random-split run discussed in \S\ref{sec:feynman30}
(superseded by 12/30) and have not been recomputed under the corrected
protocol; do not cite this paragraph's numbers without doing so first.
Reported statistics use the Kaggle run.
Independent replication of the $p=0.00020$ Mann-Whitney result is expected to hold
under any configuration that allows multi-process PySR execution.

\textbf{Superseded by CI (2026-07-27/28).} The three-machine comparison above
is no longer the operative reproducibility story: the benchmark now runs on a
single deterministic GitHub Actions runner (\texttt{ubuntu-24.04},
\texttt{PySR\_SEED=42}) in \texttt{sednabcn/LLM-HypatiaX-REPRO}
(\texttt{exp2\_feynman}: run~30258093290, commit~\texttt{d5adf4e};
\texttt{exp2\_feynman\_extrap}: run~30329546386, commit~\texttt{c694d5f}).
This is disclosed as a limitation, not a new headline result: the extrap
job's own output still shows \texttt{null} extrapolation $\Rsq$ for most
methods on the item~4--6 equations, indicating the fixes are not yet
reflected in the commit CI is building from. Do not cite a recomputed
12/30 or 13/30 from this CI pipeline until a rerun confirms the patches
are live.

% =============================================================================
%% ── Bibliography ─────────────────────────────────────────────────────────────
%% references.bib must include entries for: sahoo2018eql, liu2022discover,
%% jin2020bayesian, koza1992genetic/koza1994genetic, raissi2019physics,
%% karniadakis2021physics, cranmer2023pysr/cranmer2023pysr,
%% hornik1989multilayer, wolpert1997no, petersen2021deep, landajuela2022unified,
%% sun2022symbolic, udrescu2020ai/udrescu2020ai, kambhampati2024llms,
%% openai2023gpt4, biggio2021neural, shojaee2023transformer, la_cava_2021_srbench,
%% cranmer2020lagrangian, wei2022chain, xu2021neural, uy2011semantically.

\bibliography{references}

% =============================================================================
\appendix

\section{Full Benchmark Case Definitions}\label{sec:benchmark_cases}
\label{app:cases}

\subsection{Automated Market Makers}\label{subsec:automated-market-makers}
\label{sec:bench_amm}

\begin{itemize}
  \item \textbf{Constant product formula}: $x \cdot y = k$.
    Tests basic AMM invariant preservation.
  \item \textbf{Spot price from AMM}: $P = y/x$.
    Algebraic; should be trivially recovered.
  \item \textbf{Spot price from AMM reserve}: $P = \Delta y / \Delta x$.
    Marginal price from reserve changes.
  \item \textbf{AMM output amount}: $\Delta y = y - k/(x + \Delta x)$.
    Tests multi-step algebraic composition.
  \item \textbf{Price slippage percentage}: $\mathrm{slip} = (P_{\mathrm{exec}} - P_{\mathrm{spot}})/P_{\mathrm{spot}}$.
  \item \textbf{Constant Product Price Impact}: quadratic impact formula.
  \item \textbf{LP share percentage}: $s = \Delta L / L_{\mathrm{total}}$.
  \item \textbf{LP fee earnings}: $e = s \cdot F \cdot V$, where $F$ is fee tier and $V$ volume.
  \item \textbf{Impermanent loss percentage}: $\mathrm{IL} = 2\sqrt{r}/(1+r) - 1$,
    where $r$ is the price ratio at withdrawal vs.\ deposit.
  \item \textbf{Impermanent loss in constant product}: related quantity for the
    constant-product setting specifically.
  \item \textbf{Impermanent loss breakeven fee rate}: the fee rate $F$ at which
    LP fees exactly offset impermanent loss.
  \item \textbf{Uniswap V3 virtual}: virtual liquidity quantity for
    concentrated positions.
  \item \textbf{Concentrated liquidity position width}: tick range formula.
  \item \textbf{Capital efficiency}: ratio of effective liquidity to deposited capital.
\end{itemize}

\subsection{Lending Protocols}\label{subsec:lending-protocols}
\label{sec:bench_lending}

\begin{itemize}
  \item \textbf{Loan-to-Value}: $\mathrm{LTV} = D/C$.
  \item \textbf{Health factor}: $h = C_{\mathrm{adj}} / D$.
  \item \textbf{Collateral ratio}: $\mathrm{CR} = C/D$.
  \item \textbf{Reserve ratio}: $\rho = R / (R + L)$.
  \item \textbf{Utilization rate of DeFi}: $u = L / (L + R)$.
  \item \textbf{Borrow APY from utilization}: kinked interest rate model.
  \item \textbf{Supply APY from borrow}: $r_s = r_b \cdot u$.
  \item \textbf{Partial liquidation amount}: amount of collateral seized.
  \item \textbf{Liquidation Price Long} and \textbf{Liquidation Price Short}:
    prices at which a position is liquidated.
  \item \textbf{Leveraged position notional}: $N = C \cdot L$, where $L$ is leverage.
  \item \textbf{Effective Leverage}: $L_{\mathrm{eff}} = N/C$.
  \item \textbf{Maximum safe leverage}: function of maintenance margin rate.
  \item \textbf{Required collateral}: collateral needed for a given position.
  \item \textbf{Cross-margin available balance}: net balance accounting for all positions.
  \item \textbf{Multi-Collateral LTV}: weighted LTV across multiple collateral types.
  \item \textbf{Liquidation price for leveraged long} and \textbf{leveraged short}:
    extended formulas accounting for funding rates.
\end{itemize}

\subsection{Derivatives Pricing}\label{subsec:derivatives-pricing}
\label{sec:bench_derivatives}

\begin{itemize}
  \item \textbf{Black--Scholes Call Price}:
    $C = S\Phi(d_1) - Ke^{-rT}\Phi(d_2)$,
    with $d_{1,2} = [\ln(S/K) + (r \pm \tfrac{1}{2}\sigma^2)T] / (\sigma\sqrt{T})$.
  \item \textbf{Black--Scholes Put Price}: $P = Ke^{-rT}\Phi(-d_2) - S\Phi(-d_1)$.
  \item \textbf{Delta}: $\Delta = \Phi(d_1)$ (call).
  \item \textbf{Gamma}: $\Gamma = \phi(d_1)/(S\sigma\sqrt{T})$.
  \item \textbf{Theta}: $\Theta = -S\phi(d_1)\sigma/(2\sqrt{T}) - rKe^{-rT}\Phi(d_2)$.
  \item \textbf{Vega}: $\mathcal{V} = S\phi(d_1)\sqrt{T}$.
  \item \textbf{Put-call parity}: $C - P = S - Ke^{-rT}$.
  \item \textbf{Forward price for derivative}: $F = Se^{rT}$.
  \item \textbf{Call option intrinsic}: $\max(S - K, 0)$.
  \item \textbf{Simple options moneyness}: $(S - K)/K$.
  \item \textbf{Convexity Adjustment}: bond convexity formula.
\end{itemize}

\subsection{Risk and Portfolio Metrics}\label{subsec:risk-and-portfolio-metrics}
\label{sec:bench_risk}

\begin{itemize}
  \item \textbf{Value at Risk at 95\,\% and 99\,\%}: parametric VaR formulas.
  \item \textbf{Expected Shortfall at 95\,\% and 99\,\%}: conditional VaR formulas.
  \item \textbf{Multi-day Value at Risk}: $\mathrm{VaR}_{T} = \mathrm{VaR}_1 \cdot \sqrt{T}$.
  \item \textbf{ES scaling for multi-day}: analogous scaling for ES.
  \item \textbf{Incremental VaR}: marginal contribution of one asset to portfolio VaR.
  \item \textbf{Portfolio VaR for two correlated assets}:
    $\mathrm{VaR}_P = \sqrt{w_1^2 \sigma_1^2 + w_2^2 \sigma_2^2 + 2w_1 w_2 \rho\sigma_1\sigma_2}$.
  \item \textbf{Correlated Portfolio VaR}: multivariate extension.
  \item \textbf{Component ES}: ES contribution decomposed by asset.
  \item \textbf{Portfolio Expected Shortfall for correlated assets}: the
    most complex risk case; requires integration over a multivariate Gaussian tail.
  \item \textbf{Annualised Portfolio tracking error}: $\mathrm{TE} = \sigma_{\varepsilon}\sqrt{252}$.
  \item \textbf{Portfolio Sharpe Ratio}: $(r_p - r_f)/\sigma_p$.
  \item \textbf{Information ratio}: $\mathrm{IR} = \alpha / \mathrm{TE}$.
  \item \textbf{Position margin ratio}: margin / position notional.
\end{itemize}

\subsection{Staking and Protocol Mechanisms}\label{subsec:staking-and-protocol-mechanisms}
\label{sec:bench_staking}

\begin{itemize}
  \item \textbf{Simple Staking APY}: $\mathrm{APY} = r$.
  \item \textbf{Compounding Staking Returns}: $(1 + r/n)^{nT} - 1$.
  \item \textbf{Staking reward for fixed lock-up}: $R = P \cdot r \cdot T$.
  \item \textbf{Validator commission adjusted}: post-commission yield.
  \item \textbf{Slashing penalty}: fractional penalty applied to staked balance.
  \item \textbf{Protocol reserve accumulation}: cumulative reserve growth formula.
  \item \textbf{Funding rate cost}: $C = N \cdot f \cdot h$, where $f$ is the funding
    rate and $h$ is the holding period in funding intervals.
  \item \textbf{APY calculation}: annualised yield from periodic rates.
\end{itemize}

\section{Benchmark Version History}\label{sec:version_history}
\label{app:versions}

\begin{table}[h]
\centering
\caption{HypatiaX benchmark version history and key changes.}
\begin{tabular}{llp{8cm}}
\toprule
Version & Cases & Key Changes \\
\midrule
v1.0 & 62 & Initial benchmark; axis-aligned splits; no trust gating. \\
v2.0 & 71 & PCA-directed splits introduced; trust gate added ($\Rsq > 0.1$). \\
v3.0 & 74 & Three hard cases added; trust gate raised to $\Rsq > 0.5$; data leakage fixed; unified executor. \\
\bottomrule
\end{tabular}
\end{table}

\section{Corrected Runtime Claim: Detailed Analysis}\label{sec:runtime_correction}
\label{app:timing}

The following table supports the analysis in the runtime note on p.\,1 and the
Contributions list (item iii).

\begin{table}[h]
\centering
\caption{Detailed timing comparison and speedup calculations (v3.0 benchmark).}
\begin{tabular}{lcccc}
\toprule
Comparison & LLM time (s) & NN time (s) & Hybrid time (s) & Speedup \\
\midrule
Mean (all 74 cases)        & 11.4 & 3.0 & 6.8 & Hybrid $2.30\times$ \emph{slower} than NN \\
Median (all 74 cases)      & 10.3 & 2.7 & 1.7 & Hybrid $1.64\times$ faster than NN \\
LLM-routed only ($n=68$)   & ---  & 2.7 & 1.56\textsuperscript{*} & Hybrid $1.73\times$ faster than NN \\
\midrule
\multicolumn{5}{l}{\textsuperscript{*}Estimated from $1.73\times$ speedup vs.\ NN median of 2.7\,s.} \\
\multicolumn{5}{l}{Previously claimed: $3.7\times$ speedup = 73\,\% reduction. \textbf{Not supported by data.}} \\
\bottomrule
\end{tabular}
\end{table}

\noindent
The source of the erroneous 73\,\% claim is unclear; it does not correspond to
any of the measured speedup values.  The most likely explanation is that an
earlier benchmark run with a different set of cases (fewer hard cases requiring
NN fallback) produced a more favourable mean speedup, and this figure was
carried forward without re-verification against v3.0 data.

\section{Outstanding Process/Credibility Items (Do Not Cite Until Resolved)}\label{sec:process_items}

Two pipeline-level issues remain open as of this revision. Neither is a
claim made in the body of this paper, but both bear on the credibility of
figures elsewhere in the reproducibility material. Full detail is
maintained in a single place to avoid the two copies drifting out of sync:
see Supplementary Material A, Section~``Remaining Open Issues''
(items~6--7), which covers (i) the project's CI dashboard (July~9 run),
which should not be cited because its own upstream notebooks report
``not found'' for that run while its summary asserts ``0 open issues'',
and (ii) the legacy exp2 baseline-lock threshold defect ($R^2 \geq 0.9999$
used in place of the pipeline's stated $R^2 \geq 0.999999$, inflating the
legacy baseline from 74/90 to 81/90), which is fixed in code but not yet
confirmed against a live CI run. Both remain open; 74/90 and not 81/90
should be treated as the current legacy figure until that confirmation
exists.

\end{document}
